\section{预解及其对换基态射的应用}\etiquette[4]{XVII.sec4}\pageoriginale

\subsection{平坦预解}\etiquette[4.1]{XVII.subsec4.1}

回顾一下, 若 $f:(X,\Aa)\rightarrow (Y,\Bb)$ 为赋环景 (或赋环拓扑斯) 的态射, 则左 $\Bb$ 模层 $\Mm$ 在 $f$ 下的逆像为左 $\Aa$ 模层 $\Aa\otimes_{f^{\ast}\Bb} f^{\ast}\Mm$.

\begin{enonce*}{命题 4.1.1}\etiquette[4.1.1]{XVII.prop4.1.1}
设 $f:(X,\Aa)\rightarrow (Y,\Bb)$ 为赋环景态射, $K$ 为右 $\Aa$ 模层, $0\rightarrow L\rightarrow M\rightarrow N\rightarrow 0$ 为左 $\Bb$ 模层的正合列. 若 $M$ 与 $N$ 平坦, 则
\begin{enumerate}
\renewcommand{\theenumi}{\roman{enumi}}
\renewcommand{\labelenumi}{(\theenumi)}
\item $L$ 平坦,

\item 下列序列正合:
\[
0\longrightarrow K\otimes_{\Aa}f^{\ast}L\longrightarrow
K\otimes_{\Aa}f^{\ast}M\longrightarrow
K\otimes_{\Aa}f^{\ast}N\longrightarrow 0{\quoi}.
\]
\end{enumerate}
\end{enonce*}

将 $f$ 分解为箭头 $(X,\Aa)\xrightarrow{u}(X,f^{\sisi{\ast}{-1}}\Bb)\xrightarrow{v}(Y,\Bb)$. 函子 $v^{\ast}$ 正合且将平坦模映为平坦模 (\sisi{IV 15}{\hyperrefext[V][V.sec1]{1}}), 所以只须对 $u$ 证明 \hyperref[XVII.prop4.1.1]{4.1.1}. 同构
$$\etiquette[4.1.1.1]{XVII.eq4.1.1.1}
K\otimes_{\Aa}(\Aa\sisi{_{f^{\ast}\Bb}\otimes}{\otimes_{f^{-1}\Bb}} P)=K\otimes_{f^{\sisi{\ast}{-1}}\Bb}P\leqno(4.1.1.1)
$$
又将问题归结为 $f$ 为恒等态射的情形.

取正合列 $0\rightarrow R\rightarrow P\rightarrow
K\rightarrow 0$, 其中 $P$ 平坦\sisi{ (\hyperref[XVII.subsubsec4.1.2]{4.1.2})}{}. 下图的各行及各列均正合:
\[
\xymatrix{
 &      &                          &  0\ar[d]                        &  &\\
 &      & R\otimes L  \ar[r]\ar[d] & P\otimes L  \ar[r]\ar[d] & K\otimes L
 \ar[r]\ar[d]& 0\\
 & 0\ar[r]     & R\otimes M  \ar[r]\ar[d] & P\otimes M  \ar[r]\ar[d] & K\otimes M
 \ar[r]\ar[d]& 0\\
 & 0\ar[r]     & R\otimes N  \ar[r] \ar[d]& P\otimes N  \ar[r]\ar[d] & K\otimes M
 \ar[r] \ar[d]& 0\\
 &      & 0         &   0   &   0   &.
}
\]

对前两列应用\pageoriginale 蛇引理, 得到
\[
0\xrightarrow{j}K\otimes L\rightarrow K\otimes M{\quoi},
\]
这就证明了 (ii). 现在去掉 $P$ 平坦的假设. 在前图中, 三列及两行为短正合列; 因而第一行正合, 即 $L$ 平坦.

\setcounter{subsubsection}{1}
\subsubsection{}\etiquette[4.1.2]{XVII.subsubsec4.1.2}
设 $\Ss$ 为赋环景范畴的一个子范畴.\nde{\MakeLinkTarget*{sga4.exp.XVII.editor-note.4.1.2.ringed-site-category}\label{XVII.ndesous}更一般地, 该理论对由赋环景组成的范畴 $\Ss$ 也成立.} $\Ss$ 中各赋环景上的模层构成范畴 $\Mm$, 它在 $\Ss$ 上阿贝尔 (\hyperref[XVII.def2.2.2]{2.2.2}), 且为纤维化及余纤维化, 其纤维是景上通常的模层范畴的对偶范畴. 以 $K(\Ss)$ 表示范畴 $K(\Mm)$ (\hyperref[XVII.exem2.2.8]{2.2.8}), 以 $K^{+}(\Ss)$ (\resp $K^{-}(\Ss)$) 表示 $K(\Ss)$ 的全子范畴, 其纤维范畴是 $S\in \Ob \Ss$ 时各范畴 $K^{+}(S)$ (\resp $K^{-}(S)$) 的对偶范畴. 范畴 $K(\Ss)$、$K^{+}(\Ss)$ 及 $K^{-}(\Ss)$ 在 $\Ss$ 上三角, 且为纤维化及余纤维化.

设\pageoriginale $f:S\rightarrow S'$ 为 $\Ss$ 中的箭头. 每个内射对象对于函子 $Rf_{\ast}$ 非循环 (\hyperref[XVII.def1.2.3]{1.2.3}), 每个平坦模层对于函子 $Lf^{\ast}$ 非循环 (\hyperref[XVII.prop4.1.1]{4.1.1} 及 \hyperref[XVII.prop1.2.6]{1.2.6}), 从而得到导出函子 (\hyperref[XVII.prop1.2.7]{1.2.7})
\begin{align*}
Rf_{\ast}:D^{+}(S) &\longrightarrow D^{+}(S')\tag{4.1.2.1}\\
Lf^{\ast}:D^{-}(S\sisi{}{'}) &\longrightarrow D^{-}(S\sisi{'}{}){\quoi}.\tag{4.1.2.2}
\end{align*}\etiquette[4.1.2.1]{XVII.eq4.1.2.1}
\etiquette[4.1.2.2]{XVII.eq4.1.2.2}%原文提示: 此段后的竖直空白过大.

若 $Rf_{\ast}$ (\resp $Lf^{\ast}$) 具有有限维数, 则在 \eqhyperref[XVII.eq4.1.2.1]{4.1.2.1} (\resp \eqhyperref[XVII.eq4.1.2.2]{4.1.2.2}) 中, 可以将上标 $+$ (\resp $-$) 换成 $b$、换成 $-$ (\resp $+$), 或将其省略 (\hyperref[XVII.prop1.2.10]{1.2.10}).\nde{\MakeLinkTarget*{sga4.exp.XVII.editor-note.4.1.2.spaltenstein-resolutions}\label{XVII.ndeKS}使用 Spaltenstein 预解, 可以无条件地省略这些上标. 交换情形可见 M.~Kashiwara 与 P.~Schapira, \emph{Categories and Sheaves}, Springer, 2010 (下文引用为 [KS]), 18.6.}

以 $D(\Ss)$ (\resp $D^{+}(\Ss)$, \resp $D^{-}(\Ss)$) 表示在 $K(\Ss)$ (\resp $K^{+}(\Ss)$, \resp $K^{-}(\Ss)$) 中将拟同构取逆所得的范畴.

\begin{enonce*}{附论 4.1.3}\etiquette[4.1.3]{XVII.sch4.1.3}
(“平凡对偶定理”).
\begin{enumerate}
\renewcommand{\theenumi}{\alph{enumi}}
\renewcommand{\labelenumi}{\theenumi)}
\item 假设所有函子 $Rf_{\ast}$ (\resp $Lf^{\ast}$) 都具有有限维数. 则:
\begin{itemize}
\item[(i)] $D(\Ss)$ 的形成与基范畴 $\Ss'\rightarrow \Ss$ 的任意变换可交换. 范畴 $D(\Ss)$ 在 $\Ss$ 上三角 \emph{(\hyperref[XVII.def2.2.3]{2.2.3})}, 且其纤维等同于 $S\in \Ob
\Ss$ 时各范畴 $D(S)$ 的对偶范畴.

\item[(ii)] 范畴 $D(\Ss)$ 在 $\Ss$ 上为余纤维化 (\resp 纤维化). 设 $f:S\rightarrow S'$; 若 $Lf^{\ast}$ (\resp $Rf_{\ast}$) 具有有限维数, 则 $D(S')$ (\resp $D(S)$) 的每个对象在 $\Ss$ 上的范畴 $D(\Ss)$ 意义下都有 $f$ 下的逆像 (\resp 正像). 逆像函子及正像函子等同于 $Lf^{\ast}$ 及 $Rf_{\ast}$.

\item[(iii)] 每个\pageoriginale $D^{-}(S')$ (\resp $D^{+}(S)$) 的对象在 $D(\Ss)$ 意义下都有逆像 (\resp 正像), 且逆像 (\resp 正像) 函子等同于 $Lf^{\ast}$ (\resp $Rf_{\ast}$).\nde{(a) 中的有限维数假设是多余的. \Cf 前面的编者注.}
\end{itemize}

\item \begin{itemize}
\item[(i)] $D^{-}(\Ss)$ (\resp $D^{+}(\Ss)$) 的形成与基范畴 $\Ss'\rightarrow \Ss$ 的任意变换可交换. 范畴 $D^{-}(\Ss)$ (\resp $D^{+}(\Ss)$) 在 $\Ss$ 上三角, 且其纤维等同于各范畴 $D^{-}(S)$ (\resp $D^{+}(S)$) 的对偶范畴.

\item[(ii)] 范畴 $D^{-}(\Ss)$ (\resp $D^{+}(\Ss)$) 在 $\Ss$ 上为纤维化 (\resp 余纤维化); “逆像”(\resp “正像”) 函子等同于导出函子 $Lf^{\ast}$ (\resp $Rf_{\ast}$).

\item[(iii)] 若 $f:S\rightarrow S'$ 使 $Rf_{\ast}$ (\resp $Lf^{\ast}$) 具有有限维数, 则 $D^{-}(S)$ (\resp $D^{+}(S')$) 的每个对象在 $D^{-}(\Ss)$ (\resp $D^{+}(\Ss)$) 意义下都有正像 (\resp 逆像), 且“正像”(\resp “逆像”) 函子等同于 $Rf_{\ast}$ (\resp $Lf^{\ast}$).
\end{itemize}
\end{enumerate}
\end{enonce*}

只须验证每种情形下均满足 \hyperref[XVII.prop2.4.1]{2.4.1} 或 \hyperref[XVII.prop2.4.3]{2.4.3} 的假设; 其他断言在作基变换 $\{S\}\rightarrow \Ss$ 或 $I\rightarrow
\Ss$ 后, 由 \hyperref[XVII.prop2.3.5]{2.3.5} 及 \hyperref[XVII.prop1.2.10]{1.2.10} 得到. 在情形 (b) 中, 平坦 (\resp 松弛) 层的复形满足 \sisi{\hyperref[XVII.prop2.4.1]{2.4.1}}{\hyperref[XVII.prop2.4.3]{2.4.3} (\resp \hyperref[XVII.prop2.4.1]{2.4.1})} 的假设. 在情形 (a) 中, 使用 \sisi{\hyperref[XVII.prop2.4.3]{2.4.3}}{\hyperref[XVII.prop2.4.1]{2.4.1} (\resp \hyperref[XVII.prop2.4.3]{2.4.3})}, 所以只须考虑那些使 $Rf_{\ast}$ (\resp $Lf^{\ast}$) 的维数被固定的 $N$ 所界的 $f$. 来证明存在“足够多”的复形 $K$, 其各分量的正像 (\resp 逆像) 对于所考虑的正像 (分别地, 逆像) 函子全部非循环. 设 $K$ 为复形,\pageoriginale $K^{\ast}$ (\resp $K_{\ast}$) 为 $K$ 的一个 Cartan--Eilenberg 预解, 由分量为松弛 (\resp 平坦) 层的复形组成.\nde{只要求对每个 $p$, 增广 $K^p\to (K^*)^{p\bigdot}$ (\resp $(K_*)^{p\bigdot}\to K^p$) 为松弛 (\resp 平坦) 预解.} 若 $n>N$, 可知 $f_{\ast}\tau''_{<n}(K^{\ast})=\tau''_{<\sisi{}{n}}f_{\ast}(K^{\ast})$ (\resp $f^{\ast}\tau''_{>-n}(K_{\ast})=\tau''_{>-n}f^{\ast}(K_{\ast})$), 其中 $\tau\sisi{^{\ast}}{''}$ 表示对于第二个微分的截断. 对 $n>2N$, 由此可知 $\tau''_{<n}(K^{\ast})$ (\resp $\tau''_{>-n}(K_{\ast})$) 满足 \sisi{\hyperref[XVII.prop2.4.3]{2.4.3}}{\hyperref[XVII.prop2.4.1]{2.4.1} (\resp \hyperref[XVII.prop2.4.3]{2.4.3})} 的假设, 并且每个复形 $K$ 都是某个拟同构的源 (\resp 靶), 而该拟同构的靶 (\resp 源) 是这样的复形.

\setcounter{subsubsection}{3}
\subsubsection{}\etiquette[4.1.4]{XVII.subsubsec4.1.4}
考虑赋环景的交换方块 \eqhyperref[XVII.eq4.1.4.1]{4.1.4.1} (\cf~\hyperref[XVII.subsec0.11]{0.11})
$$\etiquette[4.1.4.1]{XVII.eq4.1.4.1}
\vcenter{
\xymatrix{
  X' \ar[r]^{g'}\ar[d]_{f'}& X\ar[d]^{f}\\
S' \ar[r]^{g} & S\rlap{.}
}}\leqno{(4.1.4.1)}
$$
可以考虑由此图表组成的景范畴的子范畴;\nde{应当考虑此图表所定义的赋环景范畴, 一般说来严格地讲它并不是赋环景范畴的子范畴. \Cf \nderef{XVII.ndesous}.} 然后利用 \sisi{\hyperref[XVII.subsubsec2.1.2]{2.1.2}}{\hyperref[XVII.prop2.1.3]{2.1.3}}, 在 \hyperref[XVII.sch4.1.3]{4.1.3} 适用时定义从 $Lg^{\ast}Rf_{\ast}$ 到 $Rf'_{\ast}L{g'}^{\ast}$ 的“换基态射”:

\begin{enonce*}{命题 4.1.5}\etiquette[4.1.5]{XVII.prop4.1.5}
\sisi{\hyperref[XVII.subsubsec2.1.2]{2.1.2}}{\hyperref[XVII.prop2.1.3]{2.1.3}} 中从 $Lg^{\ast}Rf_{\ast}$ 到 $Rf'_{\ast}L{g'}^{\ast}$ 的换基态射在以下每种情形下都有定义:
\begin{enumerate}
\renewcommand{\theenumi}{\roman{enumi}}
\renewcommand{\labelenumi}{(\theenumi)}
\item 在范畴 $D^{+}$ 中工作, 且 $Lg^{\ast}$ 与 $L{g'}^{\ast}$ 具有有限维数.

\item 在范畴 $D^{-}$ 中工作, 且 $Rf_{\ast}$ 与 $Rf'_{\ast}$ 具有有限维数.

\item 在无界导出范畴中工作, 且 $Rf_{\ast}$、$Rf'_{\ast}$、$Lg^{\ast}$ 及 $L{g'}^{\ast}$ 具有有限维数, $Lf^{\ast}$ 与 $L{f'}^{\ast}$ (或者: $Rg_{\ast}$ 与 $Rg'_{\ast}$) 也如此.\nde{(iii) 中的有限维数假设是多余的. \Cf \nderef{XVII.ndeKS}.}
\end{enumerate}
\end{enonce*}

应当\pageoriginale 注意, (iii) 中的最后一个假设颇像一个恶作剧. 除了所考虑的景具有足够多点的情形 (\hyperref[XVII.subsubsec4.2.12]{4.2.12}), 编者不知道怎样去掉它. 这一换基态射的函子性质立即由 \sisi{\hyperref[XVII.subsubsec2.1.2]{2.1.2}}{\hyperref[XVII.prop2.1.3]{2.1.3}} 中的描述得到.

\setcounter{subsubsection}{5}
\subsubsection{}\etiquette[4.1.6]{XVII.subsubsec4.1.6}
设 $(S,\Aa)$ 为赋环景, 考虑“张量积”函子
\[

\otimes:\Mod\sisi{_d}{}(S,\Aa^{\circ})\times
\Mod\sisi{_g}{}(S,\Aa)\longrightarrow \Mod(S,C)
\]%源的第一个分量取右 Aa 模, 或左 Aa-circ 模.
其中 $C$ 为 $\Aa$ 的中心.

\hyperref[XVII.subsubsec1.1.6]{1.1.6} 的方法 (应用此方法时, 与双复形关联的单复形用直和定义) 将此函子延拓为仍记作 $\otimes$ 的函子:
\[
\otimes:K(S,\Aa^{\circ})\times K(S,\Aa)\longrightarrow
K(S,C){\quoi}.
\]

\begin{enonce*}{命题 4.1.7}\etiquette[4.1.7]{XVII.prop4.1.7}
采用上述记号:
\begin{enumerate}
\renewcommand{\theenumi}{\roman{enumi}}
\renewcommand{\labelenumi}{(\theenumi)}
\item 张量积函子可左导出 (\hyperref[XVII.def1.2.1]{1.2.1}) 为“全张量积”函子
\[
{}\overset{L}{\otimes}{}:D^{-}(S,\Aa^{\circ})\times
D^{-}(S,\Aa)\longrightarrow D^{-}(S,C){\quoi}\text{;}
\]

\item 一对上有界复形 $(K,L)$ 只要 $K$ 或 $L$ 的分量平坦, 就适合于 $\overset{L}{\otimes}$.
\end{enumerate}
\end{enonce*}

证明是标准的.

\begin{enonce*}{记号 4.1.8}\etiquette[4.1.8]{XVII.not4.1.8}
若\pageoriginale $K\in D^{-}(S,\Aa^{\circ})$ 且 $L\in D^{-}(S,\Aa)$, 以 $K\overset{L}{\otimes}L$ 表示 $K$ 与 $L$ 的全张量积. 层 $H_{k}(K\overset{L}{\otimes}L)=H^{-k}(K\overset{L}{\otimes}L)$ 记作 $\Tor_k(K,L)$, 称为 $K$ 与 $L$ 的第 $k$ 个局部超 Tor (local hypertor).
\end{enonce*}

设 $K\in D^{-}(S,\Aa^{\circ})$ 及 $L\in D^{-}(S,\Aa)$. 若对 $i>k$ 有 $H^{i}(K)=0$, 对 $i>\ell$ 有 $H^{i}(L)=0$, 可以验证, 对 $i>k+\ell$ 有 $H^{i}(K\overset{L}{\otimes}L)=0$, 并且 $H^{k+\ell}(K\overset{L}{\otimes}L)=H^{k}(K)\otimes H^{\ell}(L)$. 固定 $K$, 将 $K\overset{L}{\otimes}L$ 看作 $L$ 的函子. 若将上同调幅的定义 \hyperref[XVII.def1.2.9]{1.2.9} 应用为 $\ZZ$ 中的区间, 则此区间的上端点具有平凡性质: 它是使 $H^{i}(K)\neq0$ 的最大整数 $i$. 下端点更为有趣:

\begin{enonce*}{定义 4.1.9}\etiquette[4.1.9]{XVII.def4.1.9}
若下述等价条件成立, 就称复形 $K\in\Ob D^{b}(S,\Aa^{\circ})$ 的 \emph{\tTor 维数 $\leq n$}:
\begin{enumerate}
\renewcommand{\theenumi}{\roman{enumi}}
\renewcommand{\labelenumi}{(\theenumi)}
\item 对每个 $\Aa$ 模 $L$, 当 $k>n$ 时有 $\Tor_k(K,L)=0$.

\item 对每个 $\Aa$ 模复形 $L$, 若对 $i<\ell$ 有 $H^i(L)=0$, 则对 $i<\ell-n$ 有 $H^i(K\overset{L}{\otimes}L)=0$.

\item 存在以 $K$ 为靶、以复形 $K'$ 为源的拟同构, 其中该复形\sisi{}{有界}, 各分量平坦, 且对 $i<-n$ 为零.
\end{enumerate}
\end{enonce*}

\setcounter{subsubsection}{9}
\subsubsection{}\etiquette[4.1.10]{XVII.subsubsec4.1.10}
有如下可能的变体:
\begin{enumerate}
\renewcommand{\theenumi}{\roman{enumi}}
\renewcommand{\labelenumi}{(\theenumi)}
\item 若 $K\in\Ob D^{b}(S,\Aa^{\circ})$ 有限 \ttor 维, 且 $L\in\Ob D(S,\Aa)$, 则函子 $\overset{L}{\otimes}$ 在 $(K,L)$ 处有定义 (\hyperref[XVII.def1.2.1]{1.2.1}).

\item 若 $K\in\Ob K^{b}(S,\Aa^{\circ})$ 的各分量具有有限 \tTor 维数, 则函子 $L\mapsto K\otimes L$ $(L\in K(S,\Aa))$ 可左导出.

\item 若每个模都具有有限 \tTor 维数,\nde{此条件是多余的. 在 $\Aa$ 交换的情形下, 见 \nderef{XVII.ndeKS} 中所引 [KS], 18.6.} 则函子 $\otimes$ 导出为函子
\[
{}\overset{L}{\otimes}{}:D(S,\Aa^{\circ})\times
D(S,\Aa)\longrightarrow D(S,\Cc){\quoi}.
\]
\end{enumerate}

\subsection{Godement 松弛预解}\etiquette[4.2]{XVII.subsec4.2}\pageoriginale

\subsubsection{}\etiquette[4.2.1]{XVII.subsubsec4.2.1}
设 $S$ 为拓扑斯, $X$ 为集合. 给定 $S$ 的一个以 $X$ 为指标的点族\sisi{}{}, 等价于给定从离散拓扑空间 $X$ 所定义的拓扑斯到 $S$ 的态射 $k:\Top(X)\rightarrow S$. 此族保守 (IV \hyperref[IV.def6.4.1]{6.4.1}) 的充要条件为: 对 $S$ 上的每个层 $F$, 从 $F$ 到 $k_{\ast}k^{\ast}F$ 的伴随箭头为单态射 (\cf IV \hyperref[IV.subsubsec6.4.0]{6.4.0}).

\begin{enonce*}{定义 4.2.2}\etiquette[4.2.2]{XVII.def4.2.2}
设 $k:\Top(X)\rightarrow S$ 为拓扑斯 $S$ 的保守点族. $S$ 上阿贝尔层 $F$ 的 \emph{Godement 预解}或\emph{典范松弛预解} (相对于 $X$), 记作 $\Cc_X^{\ast}(F)$ (或简记 $\Cc^{\ast}(F)$), 是 $F$ 的如下右预解:
\begin{enumerate}
\renewcommand{\labelenumi}{\theenumi)}
\item $\Cc^{0}(F)=k_{\ast}k^{\ast}F$, 且 $\epsilon:F\rightarrow\Cc^{0}(F)$ 为伴随箭头.

\item 置 $d^{-1}=\epsilon$; 对 $n\geq0$, 递归定义

$\Cc^{n+1}(F)=\Cc^{0}(\coker d^{n-1})$, 并将 $d^n$ 定义为复合箭头
\[
d^{n}:\Cc^n(F)\rightarrow
\coker(d^{n-1})\rightarrow\Cc^{0}(\coker(d^{n-1})){\quoi}.
\]
\end{enumerate}
\end{enonce*}

\begin{enonce*}{命题 4.2.3}\etiquette[4.2.3]{XVII.prop4.2.3}
在 \emph{\hyperref[XVII.def4.2.2]{4.2.2}} 的假设下,
\begin{enumerate}
\renewcommand{\theenumi}{\roman{enumi}}
\renewcommand{\labelenumi}{(\theenumi)}
\item $\Cc^n(F)$ 为松弛层 \emph{\sisi{(V 3.7)}{\hyperrefext[V][V.def4.1]{4.1}}},

\item $\Cc^n(F)$ 是 $F$ 的正合函子,

\item 复形 $\Cc^{\ast}(F)$ 在点 $x\in X$ 处的茎是 $F_x$ 的一个典范分裂的预解.
\end{enumerate}
\end{enonce*}

层\pageoriginale $\Cc^n(F)$ 松弛, 因为它们是离散拓扑空间 $X$ 上的层 (自动松弛) 的正像.

函子 $k^{\ast}$ 与 $k_{\ast}$ 正合. 因而函子 $\Cc^0$ 正合. 对 $n\geq0$ 作归纳, 证明函子 $\mathcal{Z}^{n}(F)=\coker(d^{n-1})$ 正合, 从而 $\Cc^{n+1}=\Cc^0\mathcal{Z}^{n}$ 也正合. 置 $\Zz^{-1}(F)=F$. 层的正合列
\[
0\longrightarrow F'\longrightarrow F\longrightarrow F''\longrightarrow0
\]
对 $n\geq0$ 定义图表
\[
\xymatrix{
        &   0   \ar[d]             &  0   \ar[d]            &  0 \ar[d]\\
0  \ar[r] & \Zz^{n-1}(F\sisi{^1}{'})  \ar[r] \ar[d]& \Zz^{n-1}(F)  \ar[r] \ar[d]& \Zz^{n-1}(F'')
\ar[r] \ar[d]& 0 \\
0  \ar[r] & \Cc^n(F')  \ar[r] \ar[d]& \Cc^n(F)  \ar[r] \ar[d]&
\Cc^n (F'')  \ar[r] \ar[d]& 0 \\
 0  \ar[r] & \Zz^{n}(F')  \ar[r]\ar[d] & \Zz^{n}(F)  \ar[r] \ar[d]& \Zz^{n}(F'')  \ar[r] \ar[d]&
0\\
        &  0               &  0             &  0      &,
}
\]
其各列正合, 并且由归纳假设, 前两行也正合. 因而第三行正合, 这就证明了 (ii).

短正合列
\[
0\longrightarrow\mathcal{Z}^{n+1}(F)\longrightarrow
\Cc^{n}(F)\longrightarrow\mathcal{Z}^{n}(F)\longrightarrow0
\]
是短正合列
$$\etiquette[4.2.3.1]{XVII.eq4.2.3.1}
0\rightarrow F\xrightarrow{\epsilon}\Cc^{0}(F)\rightarrow
\mathcal{Z}^{0}(F)\rightarrow0\leqno(4.2.3.1)
$$
的特殊情形 (作代换\pageoriginale $F\mapsto\mathcal{Z}^{n-1}(F)$). 为证明 (iii), 只须证明 \eqhyperref[XVII.eq4.2.3.1]{4.2.3.1} 在 $k$ 下的逆像序列典范地分裂. 事实上, 复合箭头
\[
k^{\ast}F\xrightarrow{k^{\ast}(\epsilon)}k^{\ast}\Cc^{0}F=k^{\ast}(k_{\ast}k^{\ast})F=(k^{\ast}k_{\ast})k^{\ast}F\longrightarrow
k^{\ast}F
\]
为恒等箭头.

\begin{enonce*}[remark]{注 4.2.4}\etiquette[4.2.4]{XVII.rem4.2.4}
若 $(x_i)_{i\in X}$ 为某个景的保守点族, 则此族对应于伴随拓扑斯的一个保守点族, 因而仍可定义典范松弛预解.
\end{enonce*}

\begin{enonce*}[remark]{注 4.2.5}\etiquette[4.2.5]{XVII.rem4.2.5}
若 $S$ 为赋环拓扑斯, 则函子 $\Cc^n$ 将模映为模.
\end{enonce*}

\begin{enonce*}[remark]{注 4.2.6}\etiquette[4.2.6]{XVII.rem4.2.6}
设 $\Ss$ 为一个由景组成的范畴. 假设对每个 $S\in\Ob\Ss$, 给定一个集合 $S^d$ 及 $S$ 的一个以 $S^d$ 为指标的保守点族 $k_S:\Top(S^d)\rightarrow S$. 假设 $S^d$ 对 $S$ 具有函子性, 且图表
\[
\xymatrix{
\Top(S^d)   \ar[r]^{f^d} \ar[d]^{k_S}& \Top(T^d)\ar[d]^{k_T}\\
S        \ar[r]^{f}  & T
}
\]
交换. 则相应的典范松弛预解“对 $S$ 具有函子性”.
\end{enonce*}

若\pageoriginale $\Ss$ 为拓扑空间范畴, 可以取 $S^d$ 为底集合.

若 $\Ss$ 为在概形 $S_0$ 上局部有限型的各概形 $S$ 的平展景所成的范畴, 且对每个 $s\in S_0$, $\bar{k(s)}$ 为 $k(s)$ 的代数闭包, 则可取 $S^d$ 为 $S$ 在某个 $\bar{k(s)}$ 中取值的点的集合.

若将代数闭包 $\bar{k(s)}$ 换成足够大的代数闭扩张, 则可将“有限型”条件换成一个基数条件.

\setcounter{subsubsection}{6}
\subsubsection{}\etiquette[4.2.7]{XVII.subsubsec4.2.7}
典范松弛预解的一个优点是, 它们使我们可以构造同时适合“张量积”型函子及“正像”型函子的复形. \sisi{这}{它们}因而给出了换基态射的一种也许更易理解的构造 (\hyperref[XVII.prop4.1.5]{4.1.5}).


\begin{enonce*}{引理 4.2.8}\etiquette[4.2.8]{XVII.lem4.2.8}
设 $f:(S,\Aa)\rightarrow(S',\Bb)$ 为赋环景态射, $K\in\Ob K(S,\Aa^{\circ})$ 与 $L\in\Ob K(S',\Bb)$ 为复形, $k:\Top(X)\rightarrow S$ 为 $S$ 的一个保守点集. 假设下述条件之一成立:
\begin{enumerate}
\renewcommand{\theenumi}{\alph{enumi}}
\renewcommand{\labelenumi}{\theenumi)}
\item $K$ 与 $L$ 的上同调上有界;

\item $\Aa$ 在 $f^{\sisi{\ast}{-1}}\Bb$ 上具有有限 \tTor 维数, 且 $K\in\Ob D^b(S,\Aa^{\circ})$ 具有有限 \tTor 维数;

\item $L\in\Ob D^b(S',\Bb)$ 具有有限 \tTor 维数.
\end{enumerate}

则\pageoriginale 函子 $(K,L)\mapsto K\otimes f^{\ast}L$ 在 $(K,L)$ 处可\sisi{}{左}导出 \emph{(\hyperref[XVII.def1.2.1]{1.2.1})}; 以 ${}\overset{\LL}{\otimes}\LL f^{\ast}$ 表示其导出函子.
\begin{enumerate}
\renewcommand{\theenumi}{\roman{enumi}}
\renewcommand{\labelenumi}{(\theenumi)}
\item 若 $x$ 为 $S$ 的一个点, 则 $(K\overset{\LL}{\otimes}\LL f^{\ast}L)_x\simeq
K_x\overset{\LL}{\otimes}_{\Bb_{f(x)}}L_{f(x)}$.

\item 对 ${}\overset{\LL}{\otimes}f^{\ast}$ \eqhyperref[XVII.def1.2.1]{1.2.1}, 对 $(K,L)$ 左适合, 当且仅当对每个 $x\in X$, 对 $(K_x,L_{f(x)})$ 左适合于函子 $\otimes_{\Bb_{f(x)}}$.

\item 在假设 $(a)$ 或 $(b)$ 下, $L$ 中的函子 $K\overset{\LL}{\otimes}f^{\ast}L$ 的上同调幅为 $\subset d$ \emph{(\hyperref[XVII.def1.2.9]{1.2.9})}, 当且仅当对每个 $x\in X$, $\Bb_{f(x)}$ 模复形 $K_x$ 的 \tTor 幅为 $\subset d$.\nde{假设 (a)、(b)、(c) 全部是多余的. \Cf \nderef{XVII.ndeKS}.}
\end{enumerate}
\end{enonce*}

此引理的证明是标准的; 它表达了所考虑函子的“逐点性质”.

\setcounter{subsubsection}{8}
\subsubsection{}\etiquette[4.2.9]{XVII.subsubsec4.2.9}

设 $L$ 为复形. $L$ 的 \emph{Godement 预解} (相对于 $X$) 是与双复形 $\sisi{C}{\Cc}^{\ast}L$ 关联的单复形 (\emph{用直和定义}) (\hyperref[XVII.subsubsec1.1.4]{1.1.4}). \emph{截断至 $n$ 阶的 Godement 预解} $\tau''_{\leq n}\sisi{C}{\Cc}^{\ast}L$, 是与由 $\sisi{C}{\Cc}^{\ast}L$ 沿 Godement 微分方向截断 (\hyperref[XVII.subsubsec1.1.15]{1.1.15}) 所得双复形关联的单复形. 根据 \hyperref[XVII.prop4.2.3]{4.2.3}~(iii), $L$ 在 $X$ 上逐点伦等价于复形 $\tau''_{\leq n}\sisi{C}{\Cc}^{\ast}L$ 及 $\sisi{C}{\Cc}^{\ast}L$; 因而对于“张量积”型函子, 这些复形“与 $L$ 一样好” (\hyperref[XVII.lem4.2.8]{4.2.8}).

\begin{enonce*}{命题 4.2.10}\etiquette[4.2.10]{XVII.prop4.2.10}
设 $(S,\Aa)$ 为赋环景, $X$ 为 $S$ 的一个保守点集, $N,M\in\NN\cup\{\infty\}$, 且 $N$ 或 $M$ 有限. 设 $A$ 为 $K(S)$ 中由满足下列条件的复形 $L$ 组成的子范畴:
\begin{enumerate}
\renewcommand{\theenumi}{\alph{enumi}}
\renewcommand{\labelenumi}{\theenumi)}
\item 若 $N=\infty$, 则对充分大的 $i$ 有 $H^i(L)=0$; 若 $M=\infty$, 则对充分小的 $i$ 有 $H^i(L)=0$.

\item 对每个\pageoriginale $x\in X$ 及每个 \tTor 维数 $\leq N$ 的 $\Aa_x$ 模 $Q$, 复形 $L_x$ 左适合于函子 $\sisi{M_x}{Q}\otimes_{\Aa_x}\ast$.

\item 对每个上同调维数 $\leq M$ 的拓扑斯态射 $f:\sisi{X^{\sim}}{S}\rightarrow E$, 复形 $L$ 右适合于函子 $f_{\ast}$.
\end{enumerate}

则 $A$ 到 $K(S)$ 的嵌入, 诱导在 $A$ 中将拟同构取逆所得的范畴与范畴 $D^{\ast}(S)$ 之间的等价, 其中若 $N<M=\infty$ 则 $\ast=+$, 若 $M<N=\infty$ 则 $\ast=-$, 若 $N,M<\infty$ 则 $\ast$ 留空.
\end{enonce*}

设 $K$ 为满足 (a) 的复形. 若 $N=\infty$, 则对充分大的 $i$, 典范箭头 $a:K_1=\sisi{\sigma}{\tau}_{\leq i}K\rightarrow K$ 为拟同构. 若 $N<\infty$, 置 $K_1=K$. 设 $K_1^{\ast}$ 为一个平坦 Cartan--Eilenberg 预解, 并置 $K_2=\tau''_{\geq\sisi{}{-}N}K_1$ (沿新微分方向的截断). 从 $K_2$ 到 $K$ 的复合箭头为拟同构, 且 $K_2$ 满足 (b).

设 $K$ 为满足 (a) 及 (b) 的复形. 若 $M=\infty$, 则 $N<\infty$; 对充分大的 $i$, $b:K\rightarrow\sisi{\sigma}{\tau}_{\geq i}K_1=K_3$ 为拟同构, \sisi{从}{且} $K_3$ 仍满足 (a) 及 (b). 若 $M<\infty$, 置 $K=K_3$. 设 $K_4=\sisi{\sigma}{\tau}''_{\leq M}\sisi{C}{\Cc}^{\ast}K_3$. 则由 \hyperref[XVII.prop4.2.3]{4.2.3}~(iii), $K_4$ 满足 \sisi{(a)、(c) 及 (b)}{(a)、(b) 及 (c)}.

因此, 连续两次应用 \cite[I 2.4.2]{sga4.exp.XVII.bib.CD} 即得到此命题.

\begin{enonce*}[remark]{变体 4.2.11}\etiquette[4.2.11]{XVII.var4.2.11}
若 $N<\infty$ (\resp 若 $M<\infty$, \resp 若 $N,M<\infty$), 同样也可以用下有界 (\resp 上有界, \resp 有界) 复形来工作.
\end{enonce*}

\setcounter{subsubsection}{11}
\subsubsection{}\etiquette[4.2.12]{XVII.subsubsec4.2.12}
当所考虑的景有足够多点时, 可以如下\pageoriginale 利用这些构造, 在比 \loccit{} 更一般的假设下定义 \hyperref[XVII.subsubsec4.1.4]{4.1.4} 的\emph{换基态射}.

设给定景的交换图表 (\cf (\hyperref[XVII.subsec0.11]{0.11}))
$$\etiquette[4.2.12.1]{XVII.eq4.2.12.1}
\vcenter{
\xymatrix{
X'          \ar[r]^{g'} \ar[d]^{f'}& X\ar[d]^{f}\\
S'          \ar[r]^{g}  & S\rlap{,}
}}\leqno{(4.2.12.1)}
$$
设 $\Aa$ 为 $S$ 上的环层, $\Aa'$ 为 $S'$ 上的环层, $L$ 为 $S'$ 上的 $(\Aa',g^{\ast}\Aa)$ 双模复形. 假设:
\begin{enumerate}
\renewcommand{\theenumi}{\roman{enumi}}
\renewcommand{\labelenumi}{(\theenumi)}
\item 景 $X$ 与 $X'$ 有保守点族 $k:Q\rightarrow X$ 与 $k':Q'\rightarrow X'$, 组成交换图表
$$\etiquette[4.2.12.2]{XVII.eq4.2.12.2}
\vcenter{
\sisi{\xymatrix{
X'         \ar[r]^{g'}  \ar[d]^{k'} & X     \ar[d]^{k}  & \\
Q'         \ar[r]^{g'_0} & Q &.
}}{\xymatrix{
Q'         \ar[r]^{g'_0}  \ar[d]^{k'} & Q     \ar[d]^{k}  & \\
X'         \ar[r]^{g'} & X \rlap{.}
}}}\leqno{(4.2.12.2)}
$$

\item 下述假设之一成立:
\begin{enumerate}
\renewcommand{\theenumi}{\alph{enumi}}
\renewcommand{\labelenumi}{\theenumi)}
\item 函子 $Rf_{\ast}$ 与 $Rf'_{\ast}$ 具有有限上同调维数, 且 $L$ 上有界. 置 $s=-$.

\item $L$ 下有界, 各分量的 \tTor 维数一致有界 (作为 $g^{\ast}\Aa$ 模). 置 $s=+$.

\item $L$ 有界, 且 (a) 与 (b) 的假设成立. 置 $s=$ 为空.
\end{enumerate}
\end{enumerate}

以\pageoriginale $L\overset{\LL}{\otimes}\LL g^{\ast}:D^{s}(S,\Aa)\rightarrow
D^{s}(S', \Aa')$ 表示从 $K^{s}(S,\Aa)$ 到 $K^{s}(S',\Aa')$ 的函子 $K\mapsto L\otimes_{g^{\ast}\Aa}g^{\ast}K$ 的左导出函子 (\hyperref[XVII.def1.2.1]{1.2.1}).

将对 $D^{s}(X,\Aa)$ 中的 $K$ 定义换基态射
$$\etiquette[4.2.12.3]{XVII.eq4.2.12.3}
\varphi:L\overset{\LL}{\otimes}\LL g^{\ast}Rf_{\ast}K\longrightarrow
Rf'_{\ast}({f'}^{\ast}L\overset{\LL}{\otimes}\LL{g'}^{\ast}K)\leqno(4.2.12.3)
$$
\nde{\MakeLinkTarget*{sga4.exp.XVII.editor-note.4.2.12.3.unbounded-adjunction}\label{XVII.ndecb}无界导出范畴中的伴随使我们能够不加假设 (i) 与 (ii) 而定义这样的态射. \Cf \nderef{XVII.ndeKS}.}

由 \hyperref[XVII.prop4.2.10]{4.2.10}, 只须对属于 \hyperref[XVII.prop4.2.10]{4.2.10} 中范畴 $A$ 的复形, 对 $K$ 具有函子性地定义它, 其中
\begin{align*}
N & \geq\sup_i \text{ ($L^{i}$ 在 $g^{\ast}$ 上的 tor\sisi{}{-}维数)}\\
M & \geq\text{$Rf_{\ast}$ 的上同调维数}.
\end{align*}
这样的复形同时适合函子 $Rf_{\ast}$ 及函子 ${f'}^{\ast}L\overset{\LL}{\otimes}_{\Aa'}\sisi{}{\LL}{g'}^{\ast}$ (\hyperref[XVII.lem4.2.8]{4.2.8}). 因而典范态射
\begin{align*}
\alpha &: f_{\ast}K\longrightarrow Rf_{\ast}K\\
\beta &:
{f'}^{\ast}L\overset{\LL}{\otimes}\LL \sisi{g}{{g'}}^{\ast}K\longrightarrow
\sisi{}{{f'}^*}L\otimes \sisi{g}{{g'}}^{\ast}K
\end{align*}
为拟同构. 所以为计算 \eqhyperref[XVII.eq4.2.12.3]{4.2.12.3} 的左端 $I$, 只须取左预解 (即拟同构)
\[
\gamma:K_{1}\longrightarrow f_{\ast}K{\quoi},
\]
其中 $K_{1}$ 适合于 $L\overset{\LL}{\otimes}\LL g^{\ast}$; 有
\[
I\simeq L\otimes \sisi{\LL}{} g^{\ast}K_{1}{\quoi}.
\]

同样,\pageoriginale 为计算 \eqhyperref[XVII.eq4.2.12.3]{4.2.12.3} 的右端 $II$, 只须取右预解
\[
\delta:{f'}^{\ast}L\otimes {g'}^{\ast}K\longrightarrow K_2{\quoi},
\]
其中 $K_{2}$ 适合于 $Rf'_{\ast}$; 则
\[
II\simeq f'_{\ast}K_{2}{\quoi}.
\]
于是将换基态射定义为复合
\[
I\simeq L\otimes g^{\ast}K_{1}\xrightarrow{\gamma}L\otimes
g^{\ast}f_{\ast}K\xrightarrow{c}f'_{\ast}({f'}^{\ast}L\otimes
{g'}^{\ast}K)\xrightarrow{\sisi{\alpha}{\delta}}f'_{\ast}K_{2}\simeq II{\quoi},
\]
其中 $c$ 为实际复形层次上的通常换基态射.

此换基态射有下述刻画性质, 原则上使其相容性性质易于验证.

\begin{enonce*}{命题 4.2.13}\etiquette[4.2.13]{XVII.prop4.2.13}
在上述假设下, 设 $K\in \Ob
C^{s}(X\sisi{'}{},\Aa)$、$K_{1}\in \Ob C^{s}(\sisi{X}{S},\Aa)$ 及 $K_{2}\in \Ob C^{s}(\sisi{S}{X'},\Aa)$. 此外, 设 $u$ 为从 $K$ 到 $K_{1}$ 的 $f$ 态射, 即 $u\in
\Hom(K_{1},f_{\ast}K)\simeq \Hom(f^{\ast}K_{1},K)$, $v$ 为从 ${f'}^{\ast}L\otimes {g'}^{\ast}K$ 到 $K_{2}$ 的态射. 由 $u$、$v$ 及通常的换基态射, 得到复形态射 $c_{u,v}:L\otimes g^{\ast}K_{1}\rightarrow
f'_{\ast}\sisi{k}{K}_2$. 以 $u'$、$v'$ 及 $c'_{uv}$ 表示复合态射
\begin{align*}
u' &: K_{1}\xrightarrow{u} f_{\ast}K\rightarrow Rf_{\ast}K\\
v' &: L\overset{\LL}{\otimes}\LL{g'}^{\ast}K\rightarrow
L\otimes {g'}^{\ast}K\xrightarrow{v} K_{2}\\
c'_{u,v} &:
L\overset{\LL}{\otimes}\LL{g'}^{\ast}K_{1}\rightarrow
L\otimes {g'}^{\ast}K_{1}
\xrightarrow{c_{uv}}f'_{\ast}K_{2}\rightarrow Rf'_{\ast}K_{2}{\quoi}.
\end{align*}
下图\pageoriginale 交换:
$$\etiquette[4.2.13.1]{XVII.eq4.2.13.1}
\vcenter{
\xymatrix{
L \overset{\LL}{\otimes} \sisi{L}{\LL}g^{\ast}Rf_{\ast}K
\ar[r]^{\varphi} & Rf'_{\ast}(f'^{\ast}L
\overset{\LL}{\otimes}\sisi{L}{\LL} g'^{\ast}K) \ar[d]^{Rf'_{\ast}(v')}&\\
L \overset{\LL}{\otimes} \sisi{L}{\LL}g^{\ast}K_1 \ar[u]_{L \overset{\LL}{\otimes} \sisi{L}{\LL}g^{\ast}(u')}
\ar[r]^{c'_{uv}} & Rf'_{\ast}K_2.
}}\leqno{(4.2.13.1)}
$$
\end{enonce*}

先假设 $K_{1}=f_{\ast}K$ 且 $K_{2}={f'}^{\ast}L\otimes {g'}^{\ast}K$, 考虑下述复形交换图表:
\[
\xymatrix{
K^{a}  \ar[r] \ar[d]& K\ar[d]\\
K^{c}  \ar[r] & K^{b}
}
\]
其中 $K^{a}$ 适合于 $\sisi{L}{}{f'}^{\ast}L\overset{\LL}{\otimes}\sisi{L}{\LL}{g'}^{\ast}$, $K^{b}$ 适合于 $Rf_{\ast}$, $K^{c}$ 同时适合于这两个函子. 为构造此图表, 先构造 $K^{a}$, 再对 $K^{a}$ 与 $K$ 应用典范松弛预解 (必要时作截断).

对此图表应用函子 $f_{\ast}$ (\resp~$L\otimes {g'}^{\ast}$), 并对图表 $(K_{1},K^{a}_{1},K^{b}_{1}, K^{c}_{1})$ (\resp $(K_{2}, K^{a}_{2},
K^{b}_{2}, K^{c}_{2})$) 中的复形取预解, 得到一个交换立方体, 其中 $\approx$ 表示拟同构:
\[
\xymatrix@C=2em{
          &         &K'_1 \ar[rdd] \ar[rrrrrr]^{\approx}&        && &         &
          & K_1\ar[rdd]\\  K_1^{a'}\ar[rru]\ar[rdd]\ar[rrrrrr]^(.75){\approx} &         &      &        && & K_1^{a}\ar[rru] \ar[rdd]&       &\\
  &         &      &K^{b'}_1 \ar[rrrrrr] ^(.25){\approx}&& &         &         &   & K^{b}_1\\
          &K_1^{c'} \ar[rru]^{\approx}\ar[rrrrrr] ^(.65){\approx}&      &        && &         &K_1^{c}\ar[rru]^{\approx} \\
}
\]
其中\pageoriginale $K'_{1}$、$K^{a'}_{1}$、$K^{b'}_{1}$ 及 $K^{c'}_{1}$ 适合于函子 $L\overset{\LL}{\otimes}\sisi{L}{\LL}g^{\ast}$ (\resp \[
\xymatrix@C=2em{
        &         &K_2 \ar[rdd] \ar[rrrrrr]^{\approx}&        && &         &
          & K'_2\ar[rdd]\\
K_2^{a}\ar[rru]\ar[rdd]^{\approx}\ar[rrrrrr]^(.75){\approx} &         &      &        && & K_2^{a'}\ar[rru] \ar[rdd]^(.65){\approx}&       &\\
  &         &      &K^{b}_2 \ar[rrrrrr] ^(.25){\approx}&& &         &         &   & K^{b'}_2\\
          &K_2^{c} \ar[rru]\ar[rrrrrr] ^(.65){\approx}&      &
        && &         &K_2^{c'}\ar[rru] & &, \\
}
\], 其中 $K'_{2}$、$K^{a'}_{2}$、$K^{b'}_{2}$ 及 $K^{c'}_{2}$ 适合于函子 $Rf'_{\ast}$). 利用通常的换基态射, 得到 $S$ 上复形的一个交换棱柱图.
%%%%%%
\newcommand{\circlenum}[1]{\bigcirc \kern -.255cm \lower
  .05em\hbox{${#1}$}}

\[
\def\objectstyle{\scriptstyle}
\xymatrix@C=-.27cm{
        &         & \scriptstyle{L\otimes g^{\ast}K'_1}\ar[rrrrrr]_*+[o][F-]{\scriptscriptstyle{4}} \ar[rdd]^(.35)*+[o][F-]{\scriptscriptstyle{2}}&        && &         &
          & \scriptstyle{L\otimes g^{\ast}K_1} \ar[rdd] \ar[rrrrr]_*+[o][F-]{\scriptscriptstyle{4}}& & & & & \scriptstyle{f'_{\ast}K_2} \ar[rrrrr]_*+[o][F-]{\scriptscriptstyle{4}}\ar[rdd] & & & && \scriptscriptstyle{f'_{\ast}K'_2} \ar[rdd]\\
\scriptstyle{L\otimes g^{\ast}K_1^{a'}} \ar[rru]\ar[rdd]\ar[rrrrrr] &         &
        &        && & \scriptstyle{L\otimes g^{\ast}K_1\sisi{}{^a}} \ar[rru] \ar[rdd]\ar[rrrrr]&       &
        &&& \scriptstyle{f'_{\ast}K^{a}_2}\ar[rdd]\ar[rru] \ar[rrrrr]&
        & & & &
        \scriptstyle{f'_{\ast}K^{a'}_2}\ar[rru]_*+[o][F-]{\scriptscriptstyle{3}}\ar[rdd]^(.7)*+[o][F-]{\scriptscriptstyle{3}}^{\quad \raisebox{1cm}{}{\scriptscriptstyle{-1}}}_(.7){\text{\rotatebox{110}{$\approx$}}}\\
  &         &      & \scriptstyle{L\otimes g^{\ast}K^{\sisi{c}{b}'}_1}  \ar[rrrrrr]
        && &         &         &   & \scriptstyle{L\otimes g^{\ast}K^{b}_2}
        \ar[rrrrr]& & & & & \scriptstyle{f'_{\ast}K^{b}_2} \ar[rrrrr]& & & & & \scriptstyle{f'_{\ast}K^{b'}_2}\\
          &\scriptstyle{L\otimes g^{\ast}K^{c'}_1}
 \ar[rrrrrr]_*+[o][F-]{\scriptscriptstyle{1}}\ar[rru]_{\text{\rotatebox{50}{$\approx$}}}^*+[o][F-]{\scriptscriptstyle{2}}^{\qquad\qquad \raisebox{-.3cm}{}{\scriptscriptstyle{-1}}}
&      &        && &         &\scriptstyle{L\otimes g^{\ast}K^{c}_2} \ar[rru] \ar[rrrrr]_*+[o][F-]{\scriptscriptstyle{1}}& & & & &\scriptstyle{f'_{\ast}K^{c}_2}\ar[rru]\ar[rrrrr]_*+[o][F-]{\scriptscriptstyle{1}}&
        & & & & \scriptstyle{f'_{\ast}K^{c'}_{2'}}\ar[rru]\\
}
\]
注意 \eqhyperref[XVII.eq4.2.13.1]{4.2.13.1} 图表中的态射等同于下列复合,\pageoriginale 即可得出结论:

对于 $\varphi$: 各态射 \ding{192} 的复合;

对于 $L\overset{\LL}{\otimes}\sisi{L}{\LL}g^{\ast}(u')$: 各态射 \ding{193} 的复合, 其中一个态射是某拟同构的逆;

对于 $Rf'_{\ast}(v')$: 各态射 \ding{194} 的复合, 其中一个态射是某拟同构的逆;

对于 $\sisi{C}{c}'_{uv}$: 各态射 \ding{195} 的复合.

在一般情形下, 置 $K'_{1}=f_{\ast}K$ 及 $K'_{2}={f'}^{\ast}L\otimes g^{\ast}K$. 于是有图表
\[
\xymatrix{
&L \overset{\LL}{\otimes} \sisi{L}{\LL}g^{\ast}Rf_{\ast}K
\ar[r]^\varphi&
Rf'_{\ast}(f'^{\ast}L \overset{\LL}{\otimes}\sisi{L}{\LL}
g'^{\ast}K) \ar[d] \ar[rd]& \\
L \overset{\LL}{\otimes} \sisi{L}{\LL}g^{\ast}K_1 \ar[r]\ar[ru]& L
\overset{\LL}{\otimes} \sisi{L}{\LL}g^{\ast}K'_1 \ar[r] \ar[u]&
Rf'_{\ast}K'_2 \ar[r]&  Rf'_{\ast}K_2,
}
\]
其中各三角形交换, 内部的矩形根据上述论证也交换. 因而外边界交换, 这就证明了 \hyperref[XVII.prop4.2.13]{4.2.13}.

\subsection{换基定理}\etiquette[4.3]{XVII.subsec4.3}\pageoriginale

将用导出范畴的语言重新表述固有态射的换基定理 \hyperrefext[XII][XII.thm5.1]{5.1} (iii).

\begin{enonce*}{定理 4.3.1}\etiquette[4.3.1]{XVII.thm4.3.1}
设
\[
\xymatrix{
X' \ar[rr]^{g'} \ar[d]^{f'}&& X\ar[d]^{f}\\
S' \ar[rr]^{g}& & S
}
\]
为概形的笛卡尔图表, 其中 $f$ 固有. 设 $\Aa$ 与 $\Aa'$ 为 $S$ 及 $S'$ 上的挠环层, $L$ 为上有界的 $(\Aa',g^{\ast}\Aa)$ 双模复形. 假设 $f$ 的各纤维维数有界. 则 \hyperref[XVII.subsubsec4.2.12]{4.2.12} 的假设 {\rm(a)} 满足, 从 $D^{-}(X,f^{\ast}\Aa)$ 到 $D^{-}(S',\Aa')$ 的函子之间的换基态射 \eqhyperref[XVII.eq4.2.12.3]{4.2.12.3}
\[
\varphi_K:L\overset{\LL}{\otimes}\sisi{L}{\LL}g^{\ast}\sisi{}{Rf_*}K\longrightarrow
Rf'_{\ast}({f'}^{\ast}L\overset{\LL}{\otimes}\sisi{L}{\LL}{g'}^{\ast}K){\quoi},
\]
为同构.
\end{enonce*}

所考虑的函子是三角的且远端有界 (way-out) \cite[I 7]{sga4.exp.XVII.bib.RD}, 所以为验证 $\varphi_{K}$ 为同构, 只须验证各 $\varphi_{ℋ^{i}(K)}$ 为同构. 因而可假设 $K$ 为集中在次数 $0$ 的挠 $f^{\ast}\Aa$ 模层.

为\pageoriginale 验证 $\varphi_{K}$ 为同构, 只须在应用从 $\Aa'$ 模范畴到 $\ZZ/n\ZZ$ 模范畴的遗忘函子后验证之, 其中 $n\Aa=n\Aa'=0$.

于是归结为 $\Aa'=\ZZ/n\ZZ$ 且 $n\Aa=0$ 的情形. 此时可以将 $L$ 换成导出范畴中与之同构的一个上有界复形, 其各分量在 $g^{\ast}\Aa$ 上平坦. 由于所考虑的函子对 $L$ 是三角的且远端有界 (way-out), 只须在将 $L$ 换成其某个分量后验证 $\varphi_{K}$ 为同构. 因而可以假设 $L$ 为置于次数 $0$ 的平坦 $g^{\ast}\Aa$ 模层. 须证明态射
\[
\varphi^{q}:L\otimes g^{\ast}R^{q}f_{\ast}K\longrightarrow
R^{q}f'_{\ast}({f'}^{\ast}L\otimes g^{\ast}K)
\]
为同构. 设 $s$ 为 $S'$ 的几何点, $X_s$ 为 $X'$ 在 $s$ 处的纤维, $K_s$ 为 $K$ 在 $X_s$ 上的逆像. 根据 \hyperrefext[XII][XII.thm5.1]{5.1} (iii), $\varphi^{q}$ 在 $s$ 处的茎为态射
\[
\varphi^{q}_s:L_{s}\otimes H^{q}(X_s,K)\longrightarrow H^{q}(X_s,
L_s\otimes K){\quoi}.
\]
根据 $D$. Lazard 的定理 \cite{sga4.exp.XVII.bib.La}, 平坦 $\Aa_s$ 模 $L_s$ 是自由 $\Aa_s$ 模 $L_{s,i}$ 的滤过正向极限; 函子 $H^{q}(X_s,\ast)$ 与滤过正向极限可交换\sisi{}{ (\hyperrefext[VII][VII.prop3.3]{3.3})}. 因而态射 $\varphi^{q}$ 为态射
\[
\varphi^{q}_i:L_{s,i}\otimes H^{q}(X_s,K)\longrightarrow H^{q}(X_s,
L_{s,i}\otimes K){\quoi}.
\]
的正向极限. 后者显然是同构; 因而 $\varphi^{q}_s$ 也为同构, $\varphi\sisi{}{^q}$ 为同构.

\begin{enonce*}[remark]{变体 4.3.2}\etiquette[4.3.2]{XVII.var4.3.2}
\begin{enumerate}
\renewcommand{\theenumi}{\roman{enumi}}
\renewcommand{\labelenumi}{(\theenumi)}
\item 当\pageoriginale $L$ 为有界复形, 且其各分量在 $g^{\ast}\Aa$ 上具有有限 \tTor 维数时, \hyperref[XVII.subsubsec4.2.12]{4.2.12} 的假设 (c) 满足, 可以在无界导出范畴中工作.\nde{不加这些假设也可以在无界导出范畴中工作. \Cf \nderef{XVII.ndeKS}.}

\item 若 $L$ 下有界, 且其各分量在 $g^{\ast}\Aa$ 上的 \tTor 维数\sisi{有限}{一致有界}, 则即使 $\Aa$ 与 $\Aa'$ 不是挠的, \hyperref[XVII.subsubsec4.2.12]{4.2.12} 的假设 (b) 仍满足; 若层 $ℋ^{i}(K)$ 为挠层, 换基态射 $\varphi_{K}$ $(K\in
D^{+}(X,f^{\ast}\Aa)$) 将为同构.
\end{enumerate}
\end{enonce*}

\section{固有支集正像函子}\etiquette[5]{XVII.sec5}

\subsection{基本构造}\etiquette[5.1]{XVII.subsec5.1}

将遵循 \hyperref[XVII.subsec3.1]{3.1} 中所列的纲领, 对可紧化态射 $f$ 定义函子 $Rf_{!}$ (\hyperref[XVII.subsec3.2]{3.2}). 附录中将说明如何处理任意有限型分离态射的情形.\nde{\Cf \nderef{XVII.ndeConrad}.}

\subsubsection{}\etiquette[5.1.1]{XVII.subsubsec5.1.1}

设 $X$ 为拓扑斯, $U$ 为 $X$ 的开子拓扑斯 (\hyperrefext[IV][IV.def8.3]{8.3}), $j$ 为从 $U$ 到 $X$ 的嵌入态射. 在 \hyperrefext[IV][IV.subsec5.2]{5.2} 中, 定义了从 $U$ 上的集合层范畴到 $X$ 上的集合层范畴的一个“\emph{空集延拓}”函子. 此函子 $j_{!}$ 左伴随于限制函子 $j^{\ast}$.

仍以 $j_{!}$ 表示从 $U$ 上的带基点集合层范畴 (\resp 群层范畴, \resp $j^{\ast}\Aa$ 模层范畴, 其中 $\Aa$ 为 $X$ 上的环层) 到\pageoriginale $X$ 上的类似范畴的函子, 它左伴随于限制函子 $j^{\ast}$ (\sisi{IV\quad}{\hyperrefext[IV][IV.prop11.3.1]{11.3.1}}). 这三个函子 $j_{!}$ 与从模范畴到群范畴、从群范畴到带基点集合范畴的遗忘函子相容\footnote[1]{但与到集合范畴的遗忘函子不相容!
}. 若 $i:F\rightarrow X$ 为 $U$ 的闭补的嵌入, 则 $j^{\ast}j_{!}$ (典范同构于) 恒等函子, $i^{\ast}j_{!}$ 为以始对象 (或终对象, 这里两者相同) 为值的常函子. 这也刻画了 $j_{!}$. 因而所考虑的三个函子 $j_{!}$ 正合, 并嵌入相应的函子 $j_{\ast}$. 称其为\emph{零延拓}函子. 公式 $(j_1j_2)^{\ast}=j^{\ast}_2 j^{\ast}_1$ 转置为传递性同构
\[
(j_1 j_2)_{!}=j_{1!}j_{2!}{\quoi}.
\]
函子 $j_{!}$ 与换基相容:

\begin{enonce*}{引理 5.1.2}\etiquette[5.1.2]{XVII.lem5.1.2}
设 $f:X\rightarrow Y$ 为拓扑斯态射, $V$ 为 $Y$ 的开子拓扑斯, $U$ 为它在 $X$ 中的逆像:
\[
\xymatrix{
U \ar[rr]^{f'}  \ar[d]^{j'}& & V\ar[d]^{j}\\
X \ar[rr]^{f}   && Y
}
\]
则存在唯一的换基同构 $f^{\ast}j_{!}\isomto j'_{!}{f'}^{\ast}$, 使图表
$$\etiquette[5.1.2.1]{XVII.eq5.1.2.1}
\vcenter{
\xymatrix{
f^{\ast}j_{!}    \ar[r]^{\sim} \ar@{^{(}->}[d] & j'_{!}f'^{\ast} \ar@{^{(}->}[d]     \\
f^{\ast}j_{\ast}  \ar[r]^{\alpha}&j'_{\ast}f'^{\ast},
}}\leqno{(5.1.2.1)}
$$
交换, 其中 $\alpha$ 为换基态射.
\end{enonce*}

\begin{enonce*}[remark]{证明}
换基态射\pageoriginale $\varphi:j^{\ast}f_{\ast}\rightarrow f'_{\ast}{j'}^{\ast}$ 显然为同构. 它转置为换基同构 $\varphi':j'_{!}{f'}^{\ast}\isomto
f^{\ast}j_{!}$. 为验证图表 \eqhyperref[XVII.eq5.1.2.1]{5.1.2.1} 交换, 只须在限制到 $U$ 后验证之, 而这显然成立.
\end{enonce*}

\setcounter{subsubsection}{2}
\subsubsection{}\etiquette[5.1.3]{XVII.subsubsec5.1.3}
在 \hyperref[XVII.subsubsec5.1.1]{5.1.1} 的假设下, 以环层 $\Aa$ 赋环于 $X$, 模上的函子 $j_{!}$ 正合, 因而平凡地下降到导出范畴, 定义一个仍记作 $j_{!}$ 的从 $D(U,j^{\ast}\Aa)$ 到 $D(X,\Aa)$ 的函子. 仍有传递性同构
$$\etiquette[5.1.3.1]{XVII.eq5.1.3.1}
(j_1 j_2)\sisi{}{!}=j_{1!}j_{2!}{\quoi}.\leqno(5.1.3.1)
$$

在导出范畴层次上, 仍有伴随公式
$$\etiquette[5.1.3.2]{XVII.eq5.1.3.2}
\Hom_{D(X)}(j_{!}K,L)\isomto \Hom_{D(U)}
(K,j^{\ast}L){\quoi}.\leqno(5.1.3.2)
$$

\subsubsection{}\etiquette[5.1.4]{XVII.subsubsec5.1.4}
设 $f:X\rightarrow Y$ 为固有态射, $\Aa$ 为 $Y$ 上的\emph{挠}环层, 以 $\Aa$ 的逆像环层赋环于 $X$. 假设 $f$ 的各纤维维数以上界整数 $n$ 为界 (当 $Y$ 拟紧时即如此). 则已知对 $X$ 上的每个 $f^{\ast}\Aa$ 模层 $F$, 当 $q>2n$ 时层 $R^{q}f_{\ast}F$ 为零 (\hyperrefext[XII][XII.corbis5.3]{5.3 bis}). 因而函子 $f_{\ast}$ 有导出函子 (\hyperref[XVII.prop1.2.10]{1.2.10})
\[
Rf_{\ast}:D(X,f^{\ast}\Aa)\longrightarrow
D(Y,\Aa){\quoi}.
\]

若\pageoriginale $f$ 为两个固有态射的复合 $gh$, 则有传递性公式 $Rf_{\ast}=Rg_{\ast}Rh_{\ast}$, 和通常一样, 它可用余纤维化范畴的语言表述 (附论 \sisi{\hyperref[XVII.subsubsec5.1.3]{5.1.3}}{\hyperref[XVII.sch4.1.3]{4.1.3}}).


\subsubsection{}\etiquette[5.1.5]{XVII.subsubsec5.1.5}
考虑概形的交换图表
$$\etiquette[5.1.5.1]{XVII.eq5.1.5.1}
\vcenter{
\xymatrix{
U        \ar@{^{(}->}[r]_{j_1} \ar[d]^{g}      & X  \ar[d]^{f}\\
V        \ar@{^{(}->}[r]_{j_2}       & Y,
}}\leqno{(5.1.5.1)}
$$
其中 $j_1$ 与 $j_2$ 为开浸入, $f$、$g$ 为固有态射, 其纤维维数以上界整数 $n$ 为界. 设 $\Aa$ 为 $Y$ 上的挠环层, 给 $U$、$V$ 及 $X$ 赋予 $\Aa$ 的适当逆像. 由 \eqhyperref[XVII.eq5.1.3.2]{5.1.3.2}, 若 $K\in D(U)$, 则定义态射
$$\etiquette[5.1.5.2]{XVII.eq5.1.5.2}
d:j_{2!}Rg_{\ast}K\longrightarrow Rf_{\ast}j_{1!}K\leqno{(5.1.5.2)}
$$
等价于定义态射
\[
d':Rg_{\ast}K\longrightarrow j^{\ast}_2 Rf_{\ast}j_{1!}K{\quoi}.
\]

纤维积 $X'=V\times {}_YX$ 组成交换图表
$$\etiquette[5.1.5.3]{XVII.eq5.1.5.3}
\vcenter{
\xymatrix{
U \ar@{^{(}->}[r]^{k} \ar[d]^{g} & X'  \ar@{^{(}->}[r]_{j'_2} \ar[d]^{f'}& X \ar[d]^{f}          &\\
V         \ar@{=}[r] & V         \ar@{^{(}->}[r]_{j_2}  & Y   \rlap{,}
}}\leqno{(5.1.5.3)}
$$
并且,\pageoriginale 显然有 $j^{\ast}_2Rf_{\ast}j_{1!}K\simeq
Rf'_{\ast}(j'_2{}^\ast j_{1!}K)\simeq Rf'_{\ast}k_{!}K$.

态射 $k$ 为固有的开浸入; 因此 $k_{!}=k_{\ast}$, 得到同构 $Rf'_{\ast}k_{!}K\simeq
Rf'_{\ast}k_{\ast}K\simeq Rg_{\ast}K$, 从而给出所求同构 $d'$ 及态射 $d$. 此外, $d$ 在 $V$ 上的限制为同构.

\begin{enonce*}{引理 5.1.6}\etiquette[5.1.6]{XVII.lem5.1.6}
态射 \eqhyperref[XVII.eq5.1.5.2]{5.1.5.2} 为同构.
\end{enonce*}

只须验证 $d$ 在 $Y-V=F$ 上的限制为同构. 设 $X''=\fprod{X}{F}{Y}$, 将固有换基定理 (\hyperref[XVII.thm4.3.1]{4.3.1}) 应用于笛卡尔方块
\[
\xymatrix{
X'' \ar[r]_{i'} \ar[d]^{f''}& X\ar[d]^{f\sisi{'}{}}\\
F   \ar[r]_{i}& Y\rlap{,}
}
\]
其中取 $\Aa'=i^{\ast}\Aa$ 及 $L=\Aa'[0]$. 得到 $i^{\ast}Rf_{\ast}j_{1!}K\simeq Rf''_{\ast}{i'}^{\ast}j_{1!}K\simeq
Rf''_{\ast}0\simeq 0$, 因而 \eqhyperref[XVII.eq5.1.5.2]{5.1.5.2} 的两端在 $V$ 之外都为零, 且 $d$ 为同构.

\setcounter{subsubsection}{6}
\subsubsection{}\etiquette[5.1.7]{XVII.subsubsec5.1.7}
设 $S$ 为以挠环层 $\Aa$ 赋环的概形. 对每个 $S$ 概形 $\varphi:X\rightarrow S$, 置 $\Aa_X=\varphi^{\ast}\Aa$. 在由拟紧、拟分离的 $S$ 概形及 $S$ 可紧化态射组成的范畴 $(S)$ 中工作 (\hyperref[XVII.subsec3.2]{3.2}). 由 \hyperref[XVII.prop3.2.3]{3.2.3}, 此范畴、开浸入子范畴 $(S,i)$ 及固有态射子范畴 $(S,p)$ 满足 \hyperref[XVII.subsubsec3.2.4]{3.2.4} 的假设.

拟对下述数据\pageoriginale 应用 \hyperref[XVII.prop3.3.2]{3.3.2} 的理论 (采用 \hyperref[XVII.subsubsec3.3.1]{3.3.1} 的记号)

\begin{itemize}
\item[$\alpha$)] 对 $S$ 中的 $\varphi:X\rightarrow
S$ 取 $F(X)=D(X,\Aa_X)$:
\begin{itemize}
\item[(i)] 对固有的 $p:X\rightarrow Y$, 取函子 $p_{\ast}$ 为函子 $Rp_{\ast}$.

\item[(ii)] 使用通常的传递性同构 $R(pq)_{\ast}=Rp_{\ast}Rq_{\ast}$.

\item[${\rm(i')}$] 对开浸入 $i:X\rightarrow Y$, 取函子 $i_{\ast}$ 为函子 $i_{!}$.

\item[${\rm(ii')}$] 使用传递性同构 \eqhyperref[XVII.eq5.1.3.1]{5.1.3.1}.

\item[${\rm(iii)}$] 对每个交换图表 \eqhyperref[XVII.eq5.1.5.1]{5.1.5.1}, 取同构 \eqhyperref[XVII.eq5.1.5.2]{5.1.5.2} 为 $d$.
\end{itemize}
\end{itemize}

\hyperref[XVII.prop3.3.2]{3.3.2} 的条件 (a)、(b)、(c)、${\rm(a')}$、${\rm(b')}$、${\rm(c')}$ 均经验证成立, 于是证明了下述定理的 (a) 部分:

\begin{enonce*}{定理 5.1.8}\etiquette[5.1.8]{XVII.thm5.1.8}
\begin{enumerate}
\renewcommand{\theenumi}{\alph{enumi}}
\renewcommand{\labelenumi}{\theenumi)}
\item 在 \hyperref[XVII.subsubsec5.1.7]{5.1.7} 定义的范畴 $(S)$ 中, 存在且本质上唯一一种方式, 对每个箭头 $f:X\rightarrow Y$ 定义函子
\[
Rf_{!}:D(X,\Aa_X)\longrightarrow D(Y,\Aa_Y){\quoi},
\]
并对每个复合态射 $f=gh:X\xrightarrow{h}Y\xrightarrow{g}Z$ 定义 $Rf_{!}$ 与 $Rg_{!}$ $Rh_{!}$ 之间的传递性同构 $c_{\sisi{f,g}{g,h}}$, 使得
\begin{enumerate}
\renewcommand{\theenumii}{\roman{enumii}}
\renewcommand{\labelenumii}{(\theenumii)}
\item 各 $c_{f,g}$ 满足“上闭链条件”
\[
(c_{f,g}\ast Rh_{!})\circ c_{fg,h}=(Rf_{!}\ast c_{g,h})\circ
c_{f,gh}{\quoi}.
\]

\item 限制于固有态射 (\resp 开浸入) 时, 得到 \hyperref[XVII.subsubsec5.1.4]{5.1.4} (\resp~\hyperref[XVII.subsubsec5.1.3]{5.1.3}) 的变性理论.

\item 对\pageoriginale 每个图表 \eqhyperref[XVII.eq5.1.5.1]{5.1.5.1}, \eqhyperref[XVII.eq5.1.5.2]{5.1.5.2} 中的同构 $d$ 为复合同构 $c_{f,j\sisi{'}{_1}} \circ c^{-1}_{j\sisi{,f'}{_2,g}}$.
\end{enumerate}

\item 函子 $Rf_{!}$ 远端有界 \emph{\cite[I 7]{sga4.exp.XVII.bib.RD}}, 且为三角函子.
\end{enumerate}
\end{enonce*}

每个函子 $Rf_{!}$ 为函子 $Rg_{\ast}$ 与函子 $j_{!}$ 的复合, 其中 $g$ 固有且其靶拟紧, 因而纤维维数有界, $j$ 为开浸入. 这些函子远端有界且为三角函子, 因而 $Rf_{!}$ 也是如此.

\setcounter{subsubsection}{8}
\paragraph{}\etiquette[5.1.8.1]{XVII.par5.1.8.1}

\hyperref[XVII.thm5.1.8]{5.1.8} (a) 的传递性同构, 常通过由它们推导的“复合函子的谱序列”使用; 对 $K\in \Ob D(X,\Aa_X)$, 此谱序列为:
$$\etiquette[5.1.8.2]{XVII.eq5.1.8.2}
E^{pq}_2=R^{p}g_! R^{q}h_{!}(K)\Longrightarrow
R^{p+q}f_{!}(K){\quoi}.\leqno{(5.1.8.2)}
$$

为得到它, 注意根据 \hyperref[XVII.thm5.1.8]{5.1.8} b) 及 \textsc{Verdier} \cite{sga4.exp.XVII.bib.V2}, 对每个 $L\in \Ob D(Y,\Aa_Y)$, 有谱序列
$$\etiquette[5.1.8.3]{XVII.eq5.1.8.3}
E^{pq}_2=R^{p} g_!(ℋ^{q}(L))\Longrightarrow
R^{p+q}g_!(L){\quoi},\leqno{(5.1.8.3)}
$$
而 \eqhyperref[XVII.eq5.1.8.2]{5.1.8.2} 是 \eqhyperref[XVII.eq5.1.8.3]{5.1.8.3} 在 $L=Rh_{!}K$ 时的特殊情形.

\begin{enonce*}{定义 5.1.9}\etiquette[5.1.9]{XVII.def5.1.9}
\begin{enumerate}
\renewcommand{\theenumi}{\roman{enumi}}
\renewcommand{\labelenumi}{(\theenumi)}
\item 称 \hyperref[XVII.thm5.1.8]{5.1.8} 中构造的函子为“\emph{\sisi{全}{全}固有支集正像函子}”. 在不致混淆时, 将省略形容词“全”.

\item 第 $q$ 个\emph{固有支集正像函子}, 记作 $R^{q}f_{!}$, 是复合函子 $\Hh^{q}\circ Rf_{!}$.

\item 若 $X$ 为可分闭域 $k$ 上的概形, 可嵌入 $k$ 上的一个固有概形, 将\pageoriginale $\Spec (k)$ 上的模层等同于通常的模, 并以记号 $H^{q}_c(X,\quad)$、$H^{q}_{!}(X,\quad)$ 之一, 或在不致混淆时简记 $H^{q}_c(\quad)$, 表示函子 $R^{q}f_{!}$. 此时称它为“第 $q$ 个\emph{固有支集超上同调群}”, 或在它应用于集中于次数 $0$ 的复形时, 称为“第 $q$ 个\emph{固有支集上同调群}”.
\end{enumerate}
\end{enonce*}

\setcounter{subsubsection}{9}
\setcounter{paragraph}{0}
\paragraph{}\etiquette[5.1.9.1]{XVII.par5.1.9.1}
记号 $Rf_{!}$ 有滥用之处, 因为\emph{函子 $Rf_{!}$ 不是将于 \hyperref[XVII.def6.1.2]{6.1.2} 中定义的函子 $f_{!}$ 的右导出函子}. 有人更喜欢记号 $R_{!}f$.

\subsubsection{}\etiquette[5.1.10]{XVII.subsubsec5.1.10}
\hyperref[XVII.thm5.1.8]{5.1.8} 证明所需的要素, 是固有换基定理及 $j$ 为固有开浸入时 $j_{\ast}$ 与 $j_{!}$ 的恒等. 由此可陈述 \hyperref[XVII.thm5.1.8]{5.1.8} 的各种变体; 其精确含义及验证留给读者.

\begin{enonce*}{变体 5.1.11}\etiquette[5.1.11]{XVII.var5.1.11}
在范畴 $(S)$ 中, 存在且本质上唯一一种方式, 对每个箭头 $f:X\rightarrow Y$ 定义从 $X$ 上的带基点集合层范畴到 $Y$ 上的带基点集合层范畴的函子 $f_{!}$, 并对每个复合态射 $f=gh$ 定义传递性同构, 满足与 \sisi{\hyperref[XVII.subsubsec5.1.4]{5.1.4}}{\hyperref[XVII.thm5.1.8]{5.1.8}} 的条件类似的条件 {\rm (i)\sisi{,}{} (ii) (iii)}.
\end{enonce*}

此函子的直接描述见 \hyperref[XVII.def6.1.2]{6.1.2}.

\begin{enonce*}{变体 5.1.12}\etiquette[5.1.12]{XVII.var5.1.12}
对\pageoriginale 每个可紧化态射 $f:X\rightarrow Y$, 可以定义从 $X$ 上的 \iInd 有限群层范畴到 $Y$ 上的带基点集合层范畴的函子 $R^{1}f_{!}$, 并且对 $f$ 的每个紧化 $f=\bar{f}j$, 有同构 $R^{1}f_{!}\simeq R^{1}\bar{f}_{\ast}\circ j_{!}$.
\end{enonce*}

这里不要求传递性公式; 唯一的问题是证明函子 $R^{1}f_{!}$ 不依赖于所选取的 $f$ 的紧化. 设 $p:(\bar{f}_1, i_{1})\rightarrow
(\bar{f}_{2}, i_{2})$ 为紧化之间的态射, 即图表
$$\etiquette[5.1.12.1]{XVII.eq5.1.12.1}
\vcenter{
\xymatrix{
X   \ar[rr]^{i_1} \ar[rrd]^{i_2} \ar[dd]_{f}& & \bar{X}_1 \ar[d]^{p}\\
    & &  \bar{X}_2 \ar[lld]^{\bar{f}_2}\\
Y   & & .
}}\leqno{(5.1.12.1)}
$$
其中 $\bar{f}_{1}=\bar{f}_2\circ p$, 且 $p$ 固有. 若 $G$ 为 $X$ 上的群\sisi{层}{层}, 则有带基点集合的正合列 (其中 $\alpha$ 为单射)
\[
0\rightarrow
R^{1}\sisi{f}{\bar{f}}_{2\ast}(p_{\ast}i_{1!}G)\xrightarrow{\alpha}R^{1}\sisi{f}{\bar{f}}_{1\ast}(i_{1!}G)\rightarrow
\bar{f}_{2\ast}R^{1}p_{\ast}(i_{1!}G){\quoi},
\]
以及从 $i_{2!}G$ 到 $p_{\ast}i_{1!}G$ 的态射, 从而经复合得到从 $R^{1}\sisi{f}{\bar{f}}_{2\ast}i_{2!}G$ 到 $R^{1}\sisi{f}{\bar{f}}_{1\ast}i_{1!}G$ 的态射 $\beta$. 与 \hyperref[XVII.par3.3.2.1]{3.3.2.1} 中一样推理, 可见问题在于证明当 $G$ 为 \iInd 有限时, 此态射总是同构.

将固有换基定理\pageoriginale 应用于 $p$ 及集合层 $i_{1!}G$, 验证得到 $p_{\ast}i_{1!}G=i_{2!}G$, 从而得到正合列
\[
0\rightarrow R^{1}
\sisi{f}{\bar{f}}_{2\ast}(i_{2!}G)\xrightarrow{\beta}R^{1}\sisi{f}{\bar{f}}_{1\ast}(i_{1!}G)\rightarrow
\bar{f}_{2\ast}R^{1}p_{\ast}i_{1!}G{\quoi}.
\]
态射 $\beta$ 为单射 (因为 $\alpha$ 如此). 将固有换基定理应用于 $p$ 及假设为 \iInd 有限的群层 $G$, 得到 $R^{1}p_{\ast}i_{1!}G=0$. 因此 $\beta$ 为满态射, 至此完成了 \hyperref[XVII.var5.1.12]{5.1.12} 的证明.

\begin{enonce*}{定义 5.1.13}\etiquette[5.1.13]{XVII.def5.1.13}
对每个景 $X$, 以 $D_{\tors}(X)$ 表示 $X$ 上的阿贝尔层范畴的导出范畴中, 由其上同调层为挠层的复形组成的全子范畴.
\end{enonce*}

\begin{enonce*}{变体 5.1.14}\etiquette[5.1.14]{XVII.var5.1.14}
在范畴 $(S)$ 中, 存在且本质上唯一一种方式, 对每个箭头 $f:X\rightarrow Y$ 定义一个三角且远端有界的函子 \emph{\cite[I 7]{sga4.exp.XVII.bib.RD}}
\[
Rf_{!}:D^{+}_{\tors}(X)\longrightarrow D^{+}_{\tors}(Y)
\]
并对每个复合态射 $f=gh$ 定义传递性同构, 满足与 \hyperref[XVII.subsubsec5.1.4]{5.1.4} 的条件类似的条件 \mbox{\rm (i) (ii) (iii)}.
\end{enonce*}

\begin{enonce*}{变体 5.1.\sisi{}{1}5}\etiquette[5.1.15]{XVII.var5.1.15}
在以挠环层赋环的 $S$ 概形范畴中工作, 其箭头为赋环概形的态射, 诱导 (\hyperref[XVII.subsec0.8]{0.8}) 一个 $S$ 可紧化的概形态射. 对每个箭头
\[
f:(X,\Aa)\longrightarrow (Y,\Bb),
\]
以\pageoriginale $Rf_{!}$ 表示从 $\Aa$ 模范畴到 $f^{\sisi{\ast}{-1}}\Bb$ 模范畴的遗忘函子与其后的通常的 $Rf_{!}$ 的复合. 这些函子有通常函子 $Rf_{!}$ 的变性性质.
\end{enonce*}

但要注意, 这里的记号 $Rf_{!}$ 仅当 $\Aa$ 在 $f^{\sisi{\ast}{-1}}\Bb$ 上有限时才合理, 取 $X=Y=\Spec(k)$ 且 $k$ 代数闭即可看出这一点.

\setcounter{subsubsection}{15}
\subsubsection{}\etiquette[5.1.16]{XVII.subsubsec5.1.16}
设 $f:X\rightarrow S$ 为可紧化态射, $j:U\rightarrow
X$ 为到 $X$ 的开浸入, $i:F\rightarrow X$ 为其闭补, $\Aa_S$ 为 $S$ 上的挠环层, $\Aa_X$ 为它在 $X$ 上的逆像, 并设 $K\in \Ob
D(X,\Aa_X)$.

复形正合列
\[
0\longrightarrow j_{!}j^{\ast}K\longrightarrow K\longrightarrow
i_{!}i^{\ast}K\longrightarrow 0
\]
由 \textsc{Verdier} \cite[\sisi{I 2.4}{II 1.1.5}]{sga4.exp.XVII.bib.CD} 在 $D(X,\Aa_X)$ 中给出一个特殊三角 (distinguished triangle). 它在 $Rf_{!}$ 下的像是特殊三角 (\hyperref[XVII.thm5.1.8]{5.1.8} (b))
$$\etiquette[5.1.16.1]{XVII.eq5.1.16.1}
R(fj)_{!}j^{\ast}K\longrightarrow Rf_{!}K\longrightarrow
R(fi)_{!}i^{\ast}K\longrightarrow R(fj)_{!}j^{\ast}K[1]\ldots\leqno{(5.1.16.1)}
$$
从而定义上同调长正合列
$$\etiquette[5.1.16.2]{XVII.eq5.1.16.2}
\ldots \sisi{}{\to}R^{q}(fj)_{!}j^{\ast}K\rightarrow
R^{q}f_{!}K\rightarrow R^{q}(fi)_{!}i^{\ast}K \xrightarrow{j} R^{q+1}(fj)_{!}j^{\ast}K\rightarrow\ldots.\leqno{(5.1.16.2)}
$$

在 $S$ 为可分闭域的谱, $K$ 为集中在次数 $0$ 的层 $𝒢$ 的特殊情形下, 序列 \eqhyperref[XVII.eq5.1.16.2]{5.1.16.2} 成为
$$\etiquette[5.1.16.3]{XVII.eq5.1.16.3}
\ldots H^{q}_c(U,𝒢)\rightarrow H^{q}_c (X,𝒢)\rightarrow
H^{q}_c(F,𝒢)\xrightarrow{j} H^{q+1}_c(U,𝒢)\ldots\leqno{(5.1.16.3)}
$$

\subsubsection{}\etiquette[5.1.17]{XVII.subsubsec5.1.17}

考虑\pageoriginale $S$ 可紧化概形的图表
\[
\xymatrix{
X \ar[rr]^{u} \ar[rd]_{f}& & Y\ar[ld]^{g}\\
  & S &,
}
\]
其中 $u$ \emph{固有}. 以 $S$ 上的一个挠环层 $\Aa_S$ 的逆像赋环于 $X$ 与 $Y$. 对 $K\in
D(Y,\Aa_Y)$, 显然的 (伴随) 态射 $K\rightarrow
Ru_{\ast}u^{\ast}K$ 在应用函子 $Rg_{!}$ 后定义态射
$$\etiquette[5.1.17.1]{XVII.eq5.1.17.1}
u^{\ast}:Rg_{!}K\longrightarrow
Rg_{!}Ru_{\ast}u^{\ast}K=Rg_{!}Ru_{!}u^{\ast}K\sisi{\eqsim}{\simeq} Rf_{!}u^{\ast}K\leqno(5.1.17.1)
$$
(\emph{固有支集上同调对于固有态射的反变性}).

当 $S$ 为代数闭域的谱, $K$ 为置于次数 $0$ 的层 $𝒢$ 时, 此态射成为
$$\etiquette[5.1.17.2]{XVII.eq5.1.17.2}
u^{\ast}:H^{q}_c(Y,𝒢)\longrightarrow H^{q}_c(X,𝒢)\leqno(5.1.17.2)
$$

\subsection{换基定理}\etiquette[5.2]{XVII.subsec5.2}

\subsubsection{}\etiquette[5.2.1]{XVII.subsubsec5.2.1}

设 $f:X\rightarrow Y$ 为拓扑斯态射, $V$ 为 $Y$ 的开集, $U=f^{-1}(V)$ 为它在 $X$ 中的逆像, $\Aa$ 为 $Y$ 上的环层, $\Aa'$ 为 $X$ 上的环层, $L$ 为上有界的 $(\Aa',
f^{\ast}\Aa)$ 双模复形:
$$\etiquette[5.2.1.1]{XVII.eq5.2.1.1}
\vcenter{
\xymatrix{
U \ar[rr]^{f'} \ar@{^{(}->}[d]^{j'}& & V \ar@{^{(}->}[d]^{j}\\
X \ar[rr]_{f} & & Y.
}}\leqno{(5.2.1.1)}
$$

于是\pageoriginale 有从 $D(V,\Aa)$ 到 $D(X,f^{\ast}\Aa):f^{\ast}j_{!}\isomto
j'_{!}{f'}^{\ast}$ 的函子之间的换基同构 \eqhyperref[XVII.eq5.1.2.1]{5.1.2.1}. 此外, 从 $D^{-}(U,{f'}^{\ast}\Aa)$ 到 $D^{-}(X,\Aa'):L\overset{\LL}{\otimes}_{f^{\ast}\Aa}j'_{!}K\simeq
j'_{!}({j'}^{\ast}L\overset{\LL}{\otimes}_{f^{\ast}\Aa}K)$ 的函子之间有一个显然的同构. 为构造它, 只须注意: 若复形 $K$ 适合于函子 ${j'}^{\ast}L\overset{\LL}{\otimes}_{f^{\ast}\Aa}$, 则复形 $j'_{!}K$ 适合于函子 $L\overset{\LL}{\otimes}_{f^{\ast}\Aa}$, 且在复形层次上存在上述类型的同构. 将这两个同构复合, 得到从 $D^{-}(V,\Aa)$ 到 $D^{-}(X,\Aa')$ 的函子之间的换基同构:
$$\etiquette[5.2.1.2]{XVII.eq5.2.1.2}
L\overset{\LL}{\otimes}_{f^{\ast}\Aa}f^{\ast}j_{!}K\isomto j'_{!}({j'}^{\ast}L\overset{\LL}{\otimes}_{{f\sisi{'}{}}^{\ast}\Aa}{f'}^{\ast}K){\quoi}.\leqno(5.2.1.2)
$$

若 $L$ 有界, 其各分量在

$f^{\ast}\Aa$ 上具有有限 \ttor 维数,\nde{此条件是多余的. \Cf \nderef{XVII.ndeKS}.} 则此同构对 $K\in\Ob D(V,\Aa)$ 仍有定义.

\subsubsection{}\etiquette[5.2.2]{XVII.subsubsec5.2.2}

设 $f:X\rightarrow S$ 与 $g:S'\rightarrow S$ 为概形态射, 并假设 $f$ 为开浸入 $j$ 及固有态射 $p$ 的复合, 后者各纤维的维数以整数 $n$ 为界 (当 $S$ 拟紧时自动成立). 这些态射组成各方块均为笛卡尔方块的图表
\[
\xymatrix{
X' \ar[rrrr]_{g'} \ar@{^{(}->}[rd]^{j'} \ar[dd]_{f'}& &         &   &
X\ar@{^{(}->}[rd]^{j} \ar[dd]_(.35){f} & \\
             &\bar{X}'\ar[rrrr]|(.75)\hole_{\bar{g}}\ar[ld]^{p'}  &&   &   & \bar{X}\ar[ld]^{p}\\
S' \ar[rrrr]_{g}  & & &&  S &.
}
\]
设\pageoriginale $\Aa$ 为 $S$ 上的挠环层, $\Aa'$ 为 $S'$ 上的挠环层, $L$ 为上有界的 $(\Aa',g^{\ast}\Aa)$ 双模复形. 以 $\Aa$ (\resp $\Aa'$) 在 $X$ 与 $\bar X$ (\resp $X'$ 与 $\bar X'$) 上的逆像 $\Aa_X$ 与 $\Aa_{\bar X}$ (\resp $\Aa'_{X'}$ 与 $\Aa'_{\bar X'}$), 给 $X$ 与 $\bar X$ (\resp $X'$ 与 $\bar X'$) 赋环. 于是, 从 $D^{-}(X,\Aa_X)$ 到 $D^{-}(\bar X',\Aa'_{\bar X'})$ 的函子 ${p'}^{\ast}L\overset{\LL}{\otimes}\bar g^{\ast}j_{!}K$ 与 $j'_{!}({f'}^{\ast}L\sisi{\otimes}{\overset{\LL}{\otimes}}
{g'}^{\ast}K)$ 之间有换基同构 \eqhyperref[XVII.eq5.2.1.2]{5.2.1.2}, 从 $D^{-}(\bar X,\Aa_{\bar X})$ 到 $D^{-}(S',\Aa')$ 的函子 $L\overset{\LL}{\otimes}g^{\ast}Rp_{\ast}K$ 与 $Rp'_{\ast}({p'}^{\ast}L\overset{\LL}{\otimes}\bar g^{\ast}K)$ 之间有换基同构 (\hyperref[XVII.thm4.3.1]{4.3.1}). 经复合, 得到从 $D^{-}(X,\Aa_X)$ 到 $D^{-}(S',\Aa')$ 的函子之间的换基同构:
$$\etiquette[5.2.2.1]{XVII.eq5.2.2.1}
L\overset{\LL}{\otimes}g^{\ast}Rp_{\ast}j_{!}K\isomto
Rp'_{\ast}j'_{!}({f'}^{\ast}L\overset{\LL}{\otimes}{g'}^{\ast}K){\quoi}.\leqno(5.2.2.1)
$$

当 $L$ 有界且其各分量具有有限 \ttor 维数时, 上述论证及 \hyperref[XVII.var4.3.2]{4.3.2}~(i) 对 $K\in\Ob D(X,\Aa_X)$ 定义了换基同构.\nde{结合前面的编者注及 \nderef{XVII.ndecb}, 可以不对 $L$ 加任何条件, 也不假设 $p$ 的各纤维维数有界, 而定义此同构.}

为引用 \hyperref[XVII.thm5.1.8]{5.1.8}, 假设 $S$ 拟紧且拟分离. 则函子 $Rp_{\ast}j_{!}=Rf_{!}$ 与 $Rp'_{\ast}j'_{!}=Rf'_{!}$ 只依赖于 $f$ 及 $f'$, 因而下述陈述有意义.

\begin{enonce*}{引理 5.2.3}\etiquette[5.2.3]{XVII.lem5.2.3}
换基同构 \eqhyperref[XVII.eq5.2.2.1]{5.2.2.1} 不依赖于所选取的 $X$ 的紧化 $f=pj$.
\end{enonce*}

设 $q:\bar X_1\rightarrow\bar X_2$ 为紧化之间的态射
$$\etiquette[5.2.3.1]{XVII.eq5.2.3.1}
\vcenter{
\xymatrix{
X \ar[rr]^{j_1} \ar[rrd]|(.65)\hole^{j_2} \ar[dd]_{f}& & \bar{X}_1 \ar[d]_{q} \ar[lldd]^{p_1}\\
  & &  \bar{X}_2 \ar[lld]^{p_2}\\
S & & .
}}\leqno{(5.2.3.1)}
$$
相对于\pageoriginale $\bar X_1$ 的换基同构 $c_1$, 是将相对于 $j_1$ 与 $p_1$ 的换基同构“复合”而得, 或等价地, 将相对于 $j_1$、$q$ 及 $p_2$ 的换基同构复合而得. 相对于 $\bar X_2$ 的同构 $c_2$, 由将相对于 $j_2$ 与 $p_2$ 的同构复合而得. 为验证 $c_1=c_2$, 只须验证相对于 $j_2$ 的同构 $c'_1$ 是相对于 $j_1$ 与 $q$ 的同构的复合 $c'_2$. 用撇号表示由各箭头 \eqhyperref[XVII.eq5.2.3.1]{5.2.3.1} 经换基 $g$ 所得的箭头; $c'_1$ 与 $c'_2$ 为同构
\[
c'_{1}, c'_{2}:L\overset{\LL}{\otimes}\bar{g}_{2}^*\ast
j_{2!}K\isomto  j'_{2!}({j'}_{2}^*\ast
L\overset{\LL}{\otimes} {g'}^{\ast}K){\quoi}.
\]
由 \eqhyperref[XVII.eq5.1.3.2]{5.1.3.2}, 为验证 $c'_1=c'_2$, 只须在限制到 $\sisi{X}{\bar X}'_2$ 后验证之, 而这立即成立.

因而, 两个紧化之间只要存在态射, 它们对应的换基同构就相等; 结论由 \hyperref[XVII.prop3.2.6]{3.2.6}~(ii) 得到.

\begin{enonce*}{引理 5.2.4}\etiquette[5.2.4]{XVII.lem5.2.4}
设 $f=uv:X\rightarrow Y\rightarrow S$ 为 $S$ 可紧化态射的复合. 相对于 $u$、$v$ 及 $f$ 的换基同构 \eqhyperref[XVII.eq5.2.2.1]{5.2.2.1} 给出交换图表 (记号同下面的 \eqhyperref[XVII.eq5.2.4.1]{5.2.4.1}):
\[
\xymatrix{
L\overset{\LL}{\otimes} g^{\ast} R\sisi{uv}{(uv)}_{!} K
\ar@<-4pt>[rr]^{\sim} \ar@{=}[d]& & R\sisi{u'v'}{(u'v')}_{!}((u'v')^{\ast}L
\overset{\LL}{\otimes}g''^{\ast}K)\ar@{=}[d]\\
L\overset{\LL}{\otimes} g^{\ast}Ru_{!}Rv_{!}K \ar@<-4pt>[r]^(.45){\sim}& Ru'_{!}(u'^{\ast}L
\overset{\LL}{\otimes}g'^{\ast}Rv_{!}\sisi{}{K}) \ar@<-4pt>[r]^(.45){\sim}& Ru'_{!}Rv'_{!}(v'^{\ast}u'^{\ast}L\overset{\LL}{\otimes}g''^{\ast}K).
}
\]
\end{enonce*}

选取\pageoriginale $(u,v)$ 的一个紧化 (\hyperref[XVII.def3.2.5]{3.2.5}):
$$\etiquette[5.2.4.1]{XVII.eq5.2.4.1}
\vcenter{
\xymatrix{
X' \ar[rr]_{g''} \ar[d]_{\sisi{u}{v}'}& & X \ar@{^{(}->}[r]^{j_1} \ar[d]_{\sisi{u}{v}}& \bar{X}
\ar[ld]_{p_1}\ar@{^{(}->}[r]^{j_2} \ar[d]& \bar{\bar{X}} \ar[ld]^{p_3}& \\
Y' \ar[rr]_{g'} \ar[d]_{\sisi{v}{u}'}& & Y \ar@{^{(}->}[r]^{j_3}\ar[d]_{\sisi{v}{u}}& \bar{Y}\ar[ld]^{p_2} &  &\\
S' \ar[rr]_{g} &  & S &      &  .
}}\leqno{(5.2.4.1)}
$$

引理 \hyperref[XVII.lem5.2.4]{5.2.4} 说相对于 $f$ 的换基同构 $c_1$ 是相对于 $u$ 与 $v$ 的换基同构的复合 $c_2$. 同构 $c_1$ 由相对于 $j_1$、$j_2$、$p_3$ 及 $p_2$ 的换基同构复合而得. 对 $c_2$, 则复合相对于 $j_1$、$p_1$、$j_3$ 及 $p_2$ 的换基同构. 因此只须比较两个同构: 相对于 $p_3$ 与 $j_2$ 的同构的复合 $c'_1$, 以及相对于 $j_3$ 与 $p_1$ 的同构的复合 $c'_2$.

考虑笛卡尔图表
\[
\xymatrix{
  \bar{X}' \ar[ddd]^{p'_1} \ar[rrrr]^{\bar{g}''} \ar@{^{(}->}[rrd]^{j'_2} &    &       &    & \bar{X}\ar@{^{(}->}[rrd]^{j_2} \ar[ddd]^{\sisi{}{p_1}}|(.38)\hole\\
                          & &\bar{\bar{X}}' \ar[rrrr]_(.3){\bar{\bar{g}}} \ar[ddd]^{p'_3}&   &&
& \bar{\bar{X}} \ar[ddd]^{\sisi{}{p_3}}\\
 & & &\\
  Y'  \ar@{^{(}->}[rrd]_{j'_3}\ar[rrrr]^(.7){g'}|(.52)\hole \ar[rdd]_{\sisi{v}{u}'} &    &&
& Y\ar[rdd]_(.65){\sisi{v}{u}} \ar@{^{(}->}[rrd]^{j_3}\\
  & &\bar{Y}' \ar[rrrr]_{\bar{g}'} \ar[ld]^{p'_2}&        &  & &\bar{Y}\ar[ld]^{p_2}\\
 &S'  \ar[rrrr]_{g} &&& & S &,
}
\]
$L'={p'_2}^{\ast}L$\sisi{}{, $K'=j_{1!}K$}\pageoriginale 以及
\[
c'_{1}, c'_{2}: L'\overset{\LL}{\otimes} {\bar{g}'}^{\ast}
R(j_{3}p_{1})_{!} K\sisi{}{'}\isomto  R(j'_{3} p'_{1})_{!} ((j'_{3}
p'_{1})^{\ast} L'\otimes {\bar{g}''}^{\ast}K\sisi{}{'}){\quoi}.
\]
由 \eqhyperref[XVII.eq5.1.3.2]{5.1.3.2}, 为验证 $c'_1=c'_2$, 只须验证 ${j'_3}^{\ast}(c'_1)={j'_3}^{\ast}(c'_2)$, 而这立即成立.

\begin{enonce*}{引理 5.2.5}\etiquette[5.2.5]{XVII.lem5.2.5}
设给定以挠环层赋环的概形的态射
\[
(S'', \Aa'')\xrightarrow{g_{2}} (S',
\Aa')\xrightarrow{g_{1}} (S,\Aa){\quoi},
\]
并设 $f:X\rightarrow S$ 为可紧化态射, 从而给出图表:
\[
\xymatrix{
X'' \ar[r]_{g'_2} \ar[d]^{f''}& X' \ar[r]_{g'_1}\ar[d]^{f'}& X \ar[d]^{f}\\
S'' \ar[r]_{g_2}& S' \ar[r]_{g_1}& S \rlap{.}
}
\]
分别以\pageoriginale $f^{\ast}\Aa$、${f'}^{\ast}\Aa'$ 及 ${f''}^{\ast}\Aa''$ 赋环于 $X$、$X'$ 及 $X''$. 则图表
\[
\xymatrix{
L(g_1g_2)^{\ast}Rf_{!}  \ar[rr]^{\sim} \ar@{=}[d]&& Rf''_{!}L(g'_1g'_2)^{\ast}\ar@{=}[d]\\
Lg_{\sisi{1}{2}}^{\ast}Lg^{\ast}_{\sisi{2}{1}} R f_{!} \ar[r]^{\sim}& Lg^{\ast}_{\sisi{1}{2}} Rf'_{!}
L\sisi{g^\ast_2}{g'^\ast_1}  \ar[r]^{\sim} & Rf''_{!} Lg'^{\ast}_{\sisi{1}{2}}  L g'^{\ast}_{\sisi{2}{1}}
}
\]
交换.
\end{enonce*}

此引理立即由函子 $j_{!}$ 以及 $p$ 固有时函子 $Rp_{\ast}$ 的类似结果得到. 含模的精确版本留给读者表述.

上述陈述归纳为下述定理.

\begin{enonce*}{定理 5.2.6}\etiquette[5.2.6]{XVII.thm5.2.6}
函子 $Rf_{!}$ 与换基可交换.
\end{enonce*}

\begin{enonce*}[remark]{变体 5.2.7}\etiquette[5.2.7]{XVII.var5.2.7}
\begin{enumerate}
\renewcommand{\theenumi}{\roman{enumi}}
\renewcommand{\labelenumi}{(\theenumi)}
\item 函子 $f_{!}$ (对于带基点集合层: \hyperref[XVII.var5.1.11]{5.1.11}) 及函子 $R^1f_{!}$ (对于群层: \hyperref[XVII.var5.1.12]{5.1.12}) 与换基可交换.

\item 函子 $Rf_{!}$ (\hyperref[XVII.var5.1.14]{5.1.14}) 与换基可交换.


\item 在 (\hyperref[XVII.subsubsec5.2.2]{5.2.2}) 的假设下, 若 $L$ 有界且各分量具有有限 \ttor 维数,\nde{如 \hyperref[XVII.subsubsec5.2.2]{5.2.2} 中所指出, 这些假设是多余的.} 则函子 $Rf_{!}:D(X)\rightarrow D(S)$ 与换基可交换.
\end{enumerate}
\end{enonce*}

下述\pageoriginale 两个命题 \hyperref[XVII.prop5.2.8]{5.2.8} 及 \hyperref[XVII.prop5.2.9]{5.2.9} 合在一起, 实质上等价于定理 \hyperref[XVII.thm5.2.6]{5.2.6}.

\begin{enonce*}{命题 5.2.8}\etiquette[5.2.8]{XVII.prop5.2.8}
设 $f:X\rightarrow S$ 为可紧化态射, $\mathcal F$ 为 $X$ 上的挠层. 对 $S$ 的每个几何点 $s$, 以 $X_s$ 表示 $X$ 在 $s$ 处的几何纤维, 以 $\mathcal F_s$ 表示所诱导的层. 则典范地有
\[
(R^{q}f_{!}(\mathcal{F}))_s=H^{q}_c(X_s,\mathcal{F}_s){\quoi}.
\]
\end{enonce*}

\begin{enonce*}{推论 5.2.8.1}\etiquette[5.2.8.1]{XVII.cor5.2.8.1}
若 $f$ 的纤维维数 $\leq d$, 则函子 $Rf_{!}$ 的上同调维数 $\leq2d$ (\hyperref[XVII.def1.2.9]{1.2.9}): 对每个挠层 $\sisi{F}{ℱ}$, 当 $i>2d$ 时有 $R^if_{!}\sisi{F}{ℱ}=0$. 因而函子 $R^{2d}f_{!}$ 在挠层范畴上右正合.
\end{enonce*}

通过换基 (\hyperref[XVII.prop5.2.8]{5.2.8}), 归结为 $S$ 为代数闭域的谱的情形. 此时若 $\bar X$ 为 $S$ 上的固有概形, 以 $X$ 为稠密 Zariski 开子概形, 且 $j:X\rightarrow\bar X$, 则 $\dim(\bar X)=\dim(X)\leq\sisi{D}{d}$. 因而对 $X$ 上的挠层 $ℱ$ 及 $i>2d$, 由 \hyperrefext[X][X.cor4.3]{4.3}、\refnoexpose{X.cor5.2},
\[
H^{i}_c(X,ℱ)=H^{i}(\bar{X},j_{!}\sisi{F}{ℱ})=0{\quoi}.
\]

\begin{enonce*}{命题 5.2.9}\etiquette[5.2.9]{XVII.prop5.2.9}
设 $f:X\rightarrow S$ 为可紧化态射, $\Aa$ 为 $S$ 上的挠环层, $K$ 为左 $f^{\ast}\Aa$ 模的上有界复形, $L$ 为右 $\Aa$ 模的上有界复形. 则典范地有
\[
L\overset{\LL}{\underset{\Aa}{\otimes}}
Rf_{!}(K)=Rf_{!}(f^{\ast}L\overset{\LL}{\underset{f^{\ast}\Aa}{\otimes}}K){\quoi}.
\]
\end{enonce*}

此\pageoriginale 命题形式地蕴含下述结果:

\begin{enonce*}{定理 5.2.10}\etiquette[5.2.10]{XVII.thm5.2.10}
设 $f:X\rightarrow S$ 为可紧化态射, $\Aa$ 为 $S$ 上的挠环层, $K\in D^{-}(X,f^{\ast}\Aa)$, 且 $d\in\ZZ$. 若 $K$ 有有限 \tTor 维数 $\leq d$ (\hyperref[XVII.def4.1.9]{4.1.9}), 则 $Rf_{!}K$ 有有限 \tTor 维数 $\leq d$.
\end{enonce*}

只须验证, 对 $S$ 上的每个右 $\Aa$ 模层 $L$,
\[
\sisi{H}{\Hh}^{-i}(L\overset{\LL}{\otimes}Rf_{!}K)=0\quad\mbox{\rm
对于}\quad i>d{\quoi}.
\]
由 \hyperref[XVII.prop5.2.9]{5.2.9}, 已知
\[
\sisi{H}{\Hh}^{-i}(L\overset{\sisi{L}{\LL}}{\otimes}Rf_{!}K)=R^{-i}f_{!}(f^{\ast}L\overset{\sisi{L}{\LL}}{\otimes}K){\quoi},
\]
且根据假设,
\[
\sisi{H}{\Hh}^{-i}(f^{\ast}L\overset{\sisi{L}{\LL}}{\otimes}K)=0\quad\mbox{\rm 对于}\quad
i>d{\quoi}.
\]
注意若 $X$ 上的模层复形 $M$ 满足对 $i>d$ 有 $\sisi{H}{\Hh}^{-i}(M)=0$, 则仍有对 $i>d$ 成立的 $\sisi{H}{\Hh}^{-i}(Rf_{!}M)=0$, 即得结论.

对于常 Noether 环层, 有下述更一般的结果:

\begin{enonce*}{定理 5.2.11}\etiquette[5.2.11]{XVII.thm5.2.11}
设 $f:X\rightarrow S$ 为景态射, $A$ 为环, $C$ 为其中心, 并设 $d\in\ZZ$. 假设
\begin{enumerate}
\renewcommand{\theenumi}{\roman{enumi}}
\renewcommand{\labelenumi}{(\theenumi)}
\item 函子 $Rf_{\ast}:D^{+}(X,C)\rightarrow D^{+}(S,C)$ 具有有限上同调维数;

\item $A$ 为右 Noether 环;

\item $S$ 有足够多点.
\end{enumerate}
对\pageoriginale 每个左 $A$ 模层的复形 $K\in\Ob D^{-}(X,A)$, 若 $K$ 有有限 \tTor 维数 $\leq d$, 则复形 $Rf_{\ast}K\in\Ob D^{-}(S,A)$ 的 \tTor 维数 $\leq d$.
\end{enonce*}

\begin{enonce*}[remark]{证明}
由 (i), 函子 $f_{\ast}$ 可导出为函子
\[
Rf_{\ast}:D(X,A)\longrightarrow D(S,A)
\]
及
\[
Rf_{\ast}:D(X,C)\longrightarrow D(S,C){\quoi}.
\]
\end{enonce*}

由 \hyperref[XVII.lem4.2.8]{4.2.8}, 只须验证对 $S$ 的每个点 $s$, $A$ 模复形 $(Rf_{\ast}K)_s$ 的 \tTor 维数 $\leq d$; 为此只须对每个有限生成右 $A$ 模 $L$ 有
\[
H^{-i}(L\overset{L}{\otimes}Rf_{\ast}K)=0\quad\mbox{\rm 对于}\quad
i>d{\quoi}.
\]

与 \hyperref[XVII.thm5.2.10]{5.2.10} 中一样推理, 可知只须“换基箭头”
$$\etiquette[5.2.11.1]{XVII.eq5.2.11.1}
L\overset{L}{\otimes}Rf_{\ast}K\longrightarrow Rf_{\ast}(L\overset{L}{\otimes}K)\leqno(5.2.11.1)
$$
为同构.

设 $L^{\ast}$ 为 $L$ 的一个由有限生成自由 $A$ 模组成的预解. 由于所考虑的函子向右远端有界, 只须在将 $L$ 换成 $L\sisi{_{\ast}}{^\ast}$ 的任意分量后, 证明 \eqhyperref[XVII.eq5.2.11.1]{5.2.11.1} 为同构; 于是归结为 $L$ 有限生成且自由的显然情形.

\begin{enonce*}[remark]{注 5.2.12}\etiquette[5.2.12]{XVII.rems5.2.12}
\begin{enumerate}
\renewcommand{\theenumi}{\roman{enumi}}
\renewcommand{\labelenumi}{(\theenumi)}
\item \hyperref[XVII.thm5.2.11]{5.2.11} 的假设\pageoriginale (iii) 很可能不必要.

\item 若省略 $A$ 的 Noether 假设, 当函子 $R^if_{\ast}$ 与滤过正向极限可交换时, \hyperref[XVII.thm5.2.11]{5.2.11} 仍成立.

\item 以\emph{常}环层赋环于 $S$ 的假设是必不可少的. 可取 $S$ 为严格局部离散赋值环的谱, $f$ 为一般点 $\eta$ 的嵌入, 环层 $\Aa$ 为在一般点等于 $\ZZ/n\ZZ$、在闭点 $s$ 等于 $\ZZ/n^2\ZZ$ 的层 ($n>1$ 与剩余特征互素), 并取所考虑的层为模 $f^{\ast}\Aa$, 即得到反例.
\end{enumerate}
\end{enonce*}

\subsection{比较定理与有限性定理}\etiquette[5.3]{XVII.subsec5.3}

\subsubsection{}\etiquette[5.3.1]{XVII.subsubsec5.3.1}
对于 $\CC$ 上的有限型概形 $X$, 可以关联其平展景 $X_{\text{\rm et}}$、超越景 $X_{\cl}$ 以及景态射 (\hyperrefext[XI][XI.sec4]{4})
\[
\epsilon:X_{\cl}\longrightarrow X_{\et}{\quoi}.
\]

若 $j:U\rightarrow X$ 为开浸入, 则交换图表
$$
\xymatrix{
U_{\cl}  \ar[r]\ar[d]   & U_{\et} \ar[d]\\
X_{\cl}  \ar[r] & X_{\et}
}
$$
将 $U^{\sim}_{\cl}$ 等同于 $X^{\sim}_{\cl}$ 的一个开子拓扑斯, 它是 $X_{\widetilde{\text{\rm et}}}$ 的开子拓扑斯 $U_{\widetilde{\text{\rm et}}}$ 的逆像.

若 $p:Y\rightarrow X$ 为固有态射, 则连续映射 $p^{\an}:Y^{\an}\rightarrow X^{\an}$ 固有且分离. 为说明这一点, 用 Chow 引理归结为 $p$ 射影的情形.

设\pageoriginale $f:X\rightarrow S$ 为 $\CC$ 上有限型概形之间的可紧化态射, $f=pj$ 为 $f$ 的一个紧化.
$$\etiquette[5.3.1.1]{XVII.eq5.3.1.1}
\vcenter{
\xymatrix{
X_{\cl} \ar[rrr]_{\epsilon} \ar[rd]^{j_{\cl}} \ar[dd]_{f_{\cl}}&  & & X \ar[rd]^{j}\ar[dd]_(.6){f}|\hole & \\
            &\bar{X}_{\cl}\ar[ld]^{p_{\cl}}  \ar[rrr]_{\epsilon} &  & &\bar{X}\ar[ld]^{p} \\
S_{\cl}  \ar[rrr]_{\epsilon} &                 &   &  S & .
}}\leqno{(5.3.1.1)}
$$
函子 $p_{\cl\ast}j_{\cl !}$ 恰为固有支集正像函子 $f_{\cl !}$ (见 \textsc{Verdier} \cite{sga4.exp.XVII.bib.V3}). 若 $S_{\cl}$, 因而 $X_{\cl}$ 分离, 则函子 $j_{\cl !}$ 将 $X_{\cl}$ 上的软层映为 $\bar X_{\cl}$ 上的软层 \cite[II 3.5.5]{sga4.exp.XVII.bib.TF}. 由于此问题在 $S$ 上是局部的, 无论如何它将 $X_{\cl}$ 上的松弛层映为 $\sisi{X}{\bar X}_{\cl}$ 上对函子 $p_{\cl\ast}$ 非循环的层. 因而 (\cf \textsc{Verdier} \cite{sga4.exp.XVII.bib.V3})
$$
Rf_{\cl !}=Rp_{\cl\ast}j_{\cl !}{\quoi}.
$$

\subsubsection{}\etiquette[5.3.2]{XVII.subsubsec5.3.2}

由 \hyperref[XVII.lem5.1.2]{5.1.2}, 在从 $D(X)$ 到 $D(\bar X_{\cl})$ 的函子之间有换基同构:
$$
\epsilon^{\ast}j_{!}\isomto j_{\cl !}\epsilon^{\ast}{\quoi}.
$$

由\pageoriginale \hyperref[XVII.subsubsec4.2.12]{4.2.12} 或 \hyperref[XVII.prop4.1.5]{4.1.5}, 在从 $D^{+}(\bar X)$ 到 $D^{+}(S_{\cl})$ 的函子之间有换基态射:
$$
\epsilon^{\ast}Rp_{\ast}\longrightarrow Rp_{\cl
\ast}\epsilon^{\ast}{\quoi}.
$$
经复合得到\emph{从} $D^{+}_{\tors}(X)$ \emph{到} $D^{+}_{\tors}(S_{\cl})$ \emph{的函子之间的换基态射} (\hyperref[XVII.def5.1.13]{5.1.13}):
$$\etiquette[5.3.2.1]{XVII.eq5.3.2.1}
\epsilon^{\ast}Rf_{!}\longrightarrow Rf_{\cl
!}\epsilon^{\ast}{\quoi}.\leqno(5.3.2.1)
$$

与 \hyperref[XVII.lem5.2.3]{5.2.3} 中一样, 可以验证它不依赖于所选取的 $f$ 的紧化; 它满足 \hyperref[XVII.lem5.2.4]{5.2.4} 类型的传递性公式.

\begin{enonce*}{定理 5.3.3}\etiquette[5.3.3]{XVII.thm5.3.3}
\emph{(比较定理).} 态射 \eqhyperref[XVII.eq5.3.2.1]{5.3.2.1} 为同构.
\end{enonce*}

根据构造, 只须处理 $f$ 固有的情形. 所考虑的函子为三角且远端有界 \cite[I 7]{sga4.exp.XVII.bib.RD}. 因而只须处理 $K$ 为置于次数 $0$ 的挠层的情形, 结果由 \hyperrefext[XVI][XVI.thm4.1]{4.1}~(i) 得到.

\begin{enonce*}[remark]{变体 5.3.4}\etiquette[5.3.4]{XVII.var5.3.4}
\begin{enumerate}
\renewcommand{\theenumi}{\roman{enumi}}
\renewcommand{\labelenumi}{(\theenumi)}
\item 若 $\Aa$ 为 $S_{\text{\rm et}}$ 上的挠环层, 同样得到从 $D(X,f^{\ast}\Aa)$ 到 $D(S,\epsilon^{\ast}\Aa)$ 的函子之间的同构 \eqhyperref[XVII.eq5.3.2.1]{5.3.2.1}.

\item 在低次时, 对带基点集合层及 \iInd 有限群层得到通常的变体.
\end{enumerate}
\end{enonce*}

\begin{enonce*}{推论 5.3.5}\etiquette[5.3.5]{XVII.cor5.3.5}
设\pageoriginale $X$ 为 $\CC$ 上的有限型分离概形, 可紧化于 $\CC$ 上的固有概形, $F$ 为 $X$ 上的挠层. 则
$$
H^{q}_{!}(X,F)\isomto H^{q}_c(X_{\cl},\epsilon^{\ast}F){\quoi}.
$$
\end{enonce*}

\begin{enonce*}{定理 5.3.6}\etiquette[5.3.6]{XVII.thm5.3.6}
设 $f:X\rightarrow S$ 为可紧化的有限表示态射, $A$ 为左 Noether 挠环. 函子 $Rf_{!}:D(X,A)\rightarrow D(S,A)$ 将具有可构造上同调的复形映为具有可构造上同调的复形.
\end{enonce*}

此断言在 $S$ 上是局部的, 可以假设其仿射. 通过取极限, 归结为 $S$ 为 Noether 概形的情形. 具有可构造上同调的复形组成 $D(S,A)$ 的三角子范畴. 由于函子 $Rf_{!}$ 远端有界, 只须验证若 $\mathcal F$ 为可构造的 $A$ 模层, 则层 $R^qf_{!}\mathcal F$ 可构造. 写成 $f=pj$, 其中 $p$ 固有, $j$ 为开浸入. 则 $R^qf_{!}\mathcal F=R^qp_{\ast}(j_{!}\mathcal F)$. 层 $j_{!}\mathcal F$ 可构造, 结论由 \hyperrefext[XIV][XIV.thm1.1]{1.1} 得到.

\setcounter{subsubsection}{6}
\subsubsection{}\etiquette[5.3.7]{XVII.subsubsec5.3.7}
在低次时, 对一个可构造的 \iInd 有限集合层或群层的固有支集正像, 有类似陈述.

\begin{enonce*}{推论 5.3.8}\etiquette[5.3.8]{XVII.cor5.3.8}
设 $X$ 为可分闭域 $k$ 上的有限型分离概形, $\Ff$ 为 $X$ 上的可构造挠层. 假设 $X$ 可紧化于 $k$ 上的一个固有概形.\nde{此假设自动满足. 见 \nderef{XVII.ndeConrad}.} 则群 $H^q_c(X,\mathcal F)$ 有限.
\end{enonce*}

\subsection{Künneth 公式}\etiquette[5.4]{XVII.subsec5.4}\pageoriginale

\subsubsection{}\etiquette[5.4.1]{XVII.subsubsec5.4.1}
设有景的交换图表 (\hyperref[XVII.subsec0.11]{0.11})
$$\etiquette[5.4.1.1]{XVII.eq5.4.1.1}
\vcenter{
\xymatrix@R=.5cm@C=1cm{
             &  Z'\ar[ld]_{p'}\ar[rd]^{q'} \ar[dd]^{h}&    \\
X' \ar[dd]_{f}&                   &  Y' \ar[dd]^{g}\\
             &       Z\ar[ld]_{p} \ar[rd]^{q}   &    \\
X \ar[rd]           &                    & Y \ar[ld]\\
             &      S             & &.
}}\leqno{(5.4.1.1)}
$$
假设 $S$ 以环层 $\Aa$ 赋环, 其中心为 $\Cc$; 仍以 $\Aa$ 与 $\Cc$ 表示这些层在 \eqhyperref[XVII.eq5.4.1.1]{5.4.1.1} 的各景上的逆像.

若 $K$ 为 $X$ 上的右 $\Aa$ 模层, $L$ 为 $Y$ 上的左 $\Aa$ 模层, 置
\[
K\mathop{\boxtimes}_{\Aa}L=p^{\ast}\sisi{}{K}
\underset{\Aa}{\otimes}q^{\ast}L
\]
(外张量积). $X'$ (\resp $Y'$) 上的 $K$ (\resp $L$) 也类似定义. 现设 $K$ 为 $X'$ 上的右 $\Aa$ 模层, $L$ 为 $Y'$ 上的左 $\Aa$ 模层. 换基箭头 (\sisi{}{\hyperref[XVII.prop2.1.3]{2.1.3}, }\sisi{\hyperrefext[XII][XII.eq4.1]{4.1}}{\neweqhyperrefext[XII][XII.eq4.1]{4.1}}) 定义
$$\etiquette[5.4.1.2]{XVII.eq5.4.1.2}
\begin{cases}
p^{\ast}f_{\ast}K\longrightarrow h_{\ast}{p'}^{\ast}K\\
q^{\ast}g_{\ast}L\longrightarrow h_{\ast}{q'}^{\ast}L{\quoi},
\end{cases}\leqno(5.4.1.2)
$$
从而\pageoriginale 取其张量积, 得到箭头
$$\etiquette[5.4.1.3]{XVII.eq5.4.1.3}
f_{\ast}K\boxtimes_{\Aa} g_{\ast}L\longrightarrow
h_{\ast}{p'}^{\ast}K\otimes_{\Aa} h_{\ast}
{q'}^{\ast}L\longrightarrow h_{\ast}(K\boxtimes_{\Aa}L){\quoi}.\leqno(5.4.1.3)
$$

假设函子
\begin{align*}
Rf_{\ast} &:D^{+}(X',\Aa^{\circ})\longrightarrow D^{+}(X,\Aa^{\circ})\\
Rg_{\ast} &:D^{+}(Y',\Aa)\longrightarrow D^{+}(Y,\Aa)\\
Rh_{\ast} &:D^{+}(Z',\Cc)\longrightarrow D^{+}(Z,\Cc)
\end{align*}
具有有限上同调维数, 从而延拓到无界导出范畴. 因而对 $K\in\Ob D^{-}(X',\Aa^{\circ})$ 与 $L\in\Ob D^{-}(Y',\Aa)$, 自然要定义 \emph{Künneth 箭头}
$$\etiquette[5.4.1.4]{XVII.eq5.4.1.4}
Rf_{\ast}K\overset{\LL}{\boxtimes}_{\Aa}
Rg_{\ast}L\longrightarrow Rh_{\ast}(K
\overset{\LL}{\boxtimes}_{\Aa}L){\quoi}.\leqno(5.4.1.4)
$$

只在所考虑的景有足够多点时定义它, 以便使用 \hyperref[XVII.prop4.2.10]{4.2.10} 及 \hyperref[XVII.lem4.2.8]{4.2.8}. 由 \hyperref[XVII.prop4.2.10]{4.2.10}, 只须对左适合于函子 $(K,L)\mapsto K\boxtimes_{\Aa}L$ 的一对 $(K,L)$ 定义它, 其中 $K$ 与 $L$ 上有界, 并右适合于函子 $f_{\ast}$ 与 $g_{\ast}$.

设 $K_1$ 为 $f_{\ast}K$ 的平坦预解, $L_1$ 为 $g_{\ast}L$ 的平坦预解, $u:K\overset{\LL}{\boxtimes}\sisi{}{_\Aa}L\rightarrow M$ 为 $K\overset{\LL}{\boxtimes}\sisi{}{_\Aa}L$ 的适合于函子 $Rh_{\ast}$ 的预解. 有态射
\begin{gather*}
Rf_{\ast}K\overset{\LL}{\boxtimes}_{\Aa}Rg_{\ast}L\simeq
K_1\boxtimes_{\Aa}L_1\longrightarrow
f_{\ast}K\boxtimes_{\Aa}g_{\ast}L\longrightarrow\\
\longrightarrow h_{\ast}(K\boxtimes_{\Aa}L)\longrightarrow
h_{\ast}M\simeq Rh_{\ast}(K\overset{\LL}{\boxtimes}_{\Aa}L)
\end{gather*}
将 \eqhyperref[XVII.eq5.4.1.4]{5.4.1.4} 定义为其复合.

此\pageoriginale 箭头有下述刻画性质, 其证明与 \hyperref[XVII.prop4.2.13]{4.2.13} 的证明完全类似:

\setcounter{paragraph}{4}
\paragraph{}\etiquette[5.4.1.5]{XVII.par5.4.1.5}
设 $K\in\Ob C^{-}(X',\Aa^{\circ})$ 及 $L\in\Ob C^{-}(Y',\Aa)$. 置 $K_1=f_{\ast}K$、$L_1=g_{\ast}L$ 及 $M_2=K\boxtimes\sisi{}{_\Aa}L$. 则下图交换:
$$
\xymatrix@R=1.25cm{
Rf_{\ast}K
\overset{\LL}{\boxtimes}_{\Aa}Rg_{\ast}
L \ar[rrr]& & & Rh_{\ast}(K
\overset{\LL}{\boxtimes}_{\Aa}L) \ar[d]\\
K_1\overset{\LL}{\boxtimes}_{\Aa}L_{\sisi{k}{1}}
\ar[r] \ar[u]&
K_1\boxtimes\sisi{}{_\Aa}L_1 \ar[r]& h_{\ast}M_2 \ar[r]& Rh_{\ast}M_2.
}
$$

若改从态射 $K_1\rightarrow f_{\ast}K$、$L_1\rightarrow g_{\ast}L$ 及 $K\boxtimes\sisi{}{_\Aa}L\rightarrow M_2$ 出发, 也有类似的交换性 (\cf \hyperref[XVII.prop4.2.13]{4.2.13}).

\hyperref[XVII.par5.4.1.5]{5.4.1.5} 中的交换性原则上使 Künneth 态射的相容性性质易于验证.

\subsubsection{}\etiquette[5.4.2]{XVII.subsubsec5.4.2}
设 \eqhyperref[XVII.eq5.4.1.1]{5.4.1.1} 为概形的交换图表, $\Aa$ 为 $S$ 上的挠层. 假设 $f$ 与 $g$ 可紧化 (\hyperref[XVII.def3.2.1]{3.2.1}), $Z=\fprod{X}{Y}{S}$, 且 $Z'=\fprod{X'}{Y'}{S}$. 将 $f$ 与 $g$ 分解为固有态射与开浸入, 得到图表
$$
\xymatrix@C=1.75cm@R=1cm{
  X'      \sisi{\ar[r]^{p'}}{} \ar[d]^{j_1} \ar@/_1.75pc/[dd]_{f}  & Z'   \ar[d]^{j_3}\sisi{}{\ar[l]_{p'}\ar[r]^{q'}}
& \sisi{Y'\ar[l]_{q'}}{} \ar[d]^{j_2} \ar@/^1.75pc/[dd]^{g}\\
\bar{X}' \sisi{\ar[r]^{\bar{p}}}{} \ar[d]^{\bar{f}}  & \bar{Z}' \ar[d]^{\bar{h}}\sisi{}{\ar[l]_{\bar{p}}\ar[r]^{\bar{q}}} & \bar{Y}' \sisi{\ar[l]_{\bar{q}}}{}\ar[d]^{\bar{g}}\\
X            &    Z \ar[l]_{p} \ar[r]^{q}  & Y
}
$$
其中\pageoriginale $\bar Z'=\fprod{\bar X'}{\bar Y'}{S}$. 对 $X'$ 与 $Y'$ 上的复形 $K$ 与 $L$, 有一个显然的同构
$$\etiquette[5.4.2.1]{XVII.eq5.4.2.1}
j_{3!}(K\boxtimes_{\Aa}L)\isomto
j_{1!}K\boxtimes_{\Aa}j_{2!}L{\quoi}.\leqno(5.4.2.1)
$$

将此同构的逆与相对于 $(\bar f,\bar g,\bar h,j_{1!}K,j_{2!}L)$ 的 Künneth 箭头 \eqhyperref[XVII.eq5.4.1.4]{5.4.1.4}“复合”, 对\emph{ $K\in\Ob D^{-}(X',\Aa)$ 及 $L\in\Ob D^{-}(Y',\Aa)$, 得到 Künneth 箭头}
$$\etiquette[5.4.2.2]{XVII.eq5.4.2.2}
Rf_{!}K\overset{\LL}{\boxtimes}\sisi{}{_\Aa}Rg_{!}L\longrightarrow
Rh_{!}(K\overset{\LL}{\boxtimes}\sisi{}{_\Aa}L){\quoi},\leqno(5.4.2.2)
$$
可以验证它与所选取的紧化无关.

\begin{enonce*}{定理 5.4.3}\etiquette[5.4.3]{XVII.thm5.4.3}
在 \hyperref[XVII.subsubsec5.4.2]{5.4.2} 的假设下, 固有支集上同调的 Künneth 箭头 \eqhyperref[XVII.eq5.4.2.2]{5.4.2.2} 为同构.
\end{enonce*}

\begin{enonce*}[remark]{证明}
\begin{enumerate}
\renewcommand{\theenumi}{\alph{enumi}}
\renewcommand{\labelenumi}{\theenumi)}
\item 根据构造, 只须在 $f$ 与 $g$ (因而 $h$) 固有时验证结果.

\item 设给定交换图表
$$\etiquette[5.4.3.1]{XVII.eq5.4.3.1}
\vcenter{
\xymatrix@C=1.75cm@R=1.25cm{
X_1 \ar[r]^{p_1} \ar[d]& S_1 \sisi{\ar[r]^{q_1}}{}\ar[d]& Y_1\ar[d]\sisi{}{\ar[l]_{q_1}}\\
X   \ar[r]^{p}& S  \sisi{\ar[r]^{q}}{} &  Y\sisi{}{\ar[l]_q}.
}}\leqno{(5.4.3.1)}
$$
并设 $Z_1=\fprod{X_1}{Y_1}{S_1}$、$X'_1=\fprod{X'}{X_1}{X}$、$Y'_1=\fprod{Y'}{Y_1}{Y}$ 及 $Z'_1=\fprod{X'_1}{Y'_1}{S_1}=\fprod{Z'}{Z\sisi{_{1}}{}}{Z_1}$, 从而有图表 \sisi{(\hyperref[XVII.subsubsec5.4.1]{5.4.1})}{\eqhyperref[XVII.eq5.4.1.1]{5.4.1.1}} 之间的态射
\[
(S_1,X_1,Y_1,Z_1,X'_1,Y'_1,Z'_1)\longrightarrow
(S,X,Y,Z,X',Y',Z')
\]
其所有箭头统记作 $C$.

将承认\pageoriginale 下述引理:

\begin{enonce*}{引理 5.4.3.2}\etiquette[5.4.3.2]{XVII.lem5.4.3.2}
下图交换:
$$
\xymatrix{
C^{\ast}(Rf_{!} K \overset{L}{\boxtimes}_{\Aa}
Rg_{!} L) \ar@{=}[r] \ar[d]& C^{\ast}Rf_{!}K \overset{L}{\boxtimes}_{\Aa}
C^{\ast} Rg_{!} L \ar[r]^{\sim}&  Rf_{1!} C^{\ast}K
\overset{L}{\boxtimes}_{\Aa} Rg_{1!} C^{\ast}
L \ar[d]\\
C^{\ast}(Rh_{!}(K \overset{L}{\boxtimes}_{\Aa}
L)) \ar[r]^{\sim}& Rh_{\sisi{}{1}!}(C^{\ast}(K
 \overset{L}{\boxtimes}_{\Aa}  L))\ar@{=}[r] &Rh_{\sisi{}{1}!}(C^{\ast}K
 \overset{L}{\boxtimes}_{\Aa}  C^*L ).
}
$$
\end{enonce*}

在此图表中, 水平箭头为换基箭头, 竖直箭头为所考虑的两个图表 \eqhyperref[XVII.eq5.4.1.1]{5.4.1.1} 的 Künneth 箭头.

为证明 \hyperref[XVII.thm5.4.3]{5.4.3}, 只须证明 Künneth 箭头在 $Z$ 的每个几何点处为同构. 由 \hyperref[XVII.lem5.4.3.2]{5.4.3.2}, 只须在任意使 $S_1$ 为代数闭域的谱且 $p_1,q_1$ 为同构的换基 \eqhyperref[XVII.eq5.4.3.1]{5.4.3.1} 后验证之.

作这样的换基后, 图表 \eqhyperref[XVII.eq5.4.1.1]{5.4.1.1} 归结为图表
$$\vcenter{\etiquette[5.4.3.3]{XVII.eq5.4.3.3}
\xymatrix@R=1.5cm@C=2cm{
    & \fprod{X}{Y}{S} \ar[ld]_{p} \ar[rd]^{q} \ar[dd]^{f\times g}&     \\
X \ar[rd]_{f}  &                   &    Y \ar[ld]^{g}\\
    &      S             &\sisi{.}{}
}}\leqno{(5.4.3.3)}
$$
而 Künneth 箭头归结为箭头 \eqhyperref[XVII.eq5.4.4.1]{5.4.4.1}
$$\etiquette[5.4.3.4]{XVII.eq5.4.3.4}
k:Rf_{!}K\overset{\LL}{\otimes}\sisi{}{_\Aa} Rg_{!} L\longrightarrow
R(f\times g)_{!}(K\overset{\LL}{\boxtimes}\sisi{}{_\Aa}L){\quoi}.\leqno(5.4.3.4)
$$

\item 设\pageoriginale \eqhyperref[XVII.eq5.4.3.3]{5.4.3.3} 为一个其中 $f$ 与 $g$ 固有的图表. 将承认下述引理:

\begin{enonce*}{引理 5.4.3.5}\etiquette[5.4.3.5]{XVII.lem5.4.3.5}
下图交换:
$$
\xymatrix@C=1.75cm{
Rf_{!}K \overset{\sisi{L}{\LL}}{\underset{\Aa}{\otimes}} Rg_{!}L
\ar[r] \ar[d]& R(f\times g)_{!}(K
\overset{\sisi{L}{\LL}}{\underset{\Aa}{\otimes}} L) \ar@{=}[d]\\
Rf_{!}(K \overset{\sisi{L}{\LL}}{\sisi{\boxtimes}{\otimes}}_{\Aa}
f^{\ast}Rg_{!}L) \ar[r]& Rf_{!} Rp_{!}(p^{\ast}K \overset{\sisi{}{\LL}}{\otimes}\sisi{}{_\Aa} q^{\ast}L).
}
$$
\end{enonce*}

在此图表中, 水平箭头为 Künneth \sisi{箭头}{箭头}及换基箭头. 竖直箭头为 $f$ 下的换基箭头及一个传递性同构. 定理 \sisi{\hyperref[XVII.subsubsec5.5.9]{5.5.9}}{\hyperref[XVII.thm5.4.3]{5.4.3}} 由此引理及 \hyperref[XVII.thm4.3.1]{4.3.1} 得到.

勤勉的读者可以验证, 即使 $f$ 与 $g$ 不一定固有, 此引理仍然成立.
\end{enumerate}
\end{enonce*}

\setcounter{subsubsection}{3}
\subsubsection{}\etiquette[5.4.4]{XVII.subsubsec5.4.4}
设 $S$ 为概形, $f_i:X'_i\rightarrow X_i$ $(i\in I)$ 为 $S$ 概形之间的可紧化态射的有限族, $X$ 为各 $X_i$ 在 $S$ 上的积, $X'$ 为各 $X'_i$ 在 $S$ 上的积, $f:X'\rightarrow X$ 为各 $f_i$ 的积. 若如 \hyperref[XVII.subsubsec5.4.1]{5.4.1} 那样以 $\Aa$ 赋环于 $S$, 并以 $\Aa$ 的逆像赋环于各 $X_i$ 与 $X'_i$, 则更一般地, 对 $K_i\in\Ob D^{-}(X'_i,\Aa)$ 有同构
$$\etiquette[5.4.4.1]{XVII.eq5.4.4.1}
\mathop{\boxtimes}^{\LL}_{i\in I} Rf_{i!}(K_i)
\isomto Rf_{!}(\mathop{\boxtimes}^{\LL}_{i\in I}K_i){\quoi},\leqno(5.4.4.1)
$$
此同构及 \eqhyperref[XVII.eq5.4.4.1]{5.4.4.1} 两端的定义都\emph{不要求}在 $I$ 上选取全序.


